\documentclass{amsart}
% For dealing with references we use the comment environment
\usepackage{verbatim}
\newenvironment{reference}{\comment}{\endcomment}
%\newenvironment{reference}{}{}
\newenvironment{slogan}{\comment}{\endcomment}
\newenvironment{history}{\comment}{\endcomment}

% For commutative diagrams we use Xy-pic
\usepackage[all]{xy}

% We use 2cell for 2-commutative diagrams.
\xyoption{2cell}
\UseAllTwocells

% We use multicol for the list of chapters between chapters
\usepackage{multicol}

% This is generally recommended for better output
\usepackage{lmodern}
\usepackage[T1]{fontenc}

% For cross-file-references

% Package for hypertext links:
\usepackage{hyperref}

% For any local file, say "hello.tex" you want to link to please

% Theorem environments.
%
\theoremstyle{plain}
\newtheorem{theorem}[subsection]{Theorem}
\newtheorem{proposition}[subsection]{Proposition}
\newtheorem{lemma}[subsection]{Lemma}

\theoremstyle{definition}
\newtheorem{definition}[subsection]{Definition}
\newtheorem{example}[subsection]{Example}
\newtheorem{exercise}[subsection]{Exercise}
\newtheorem{situation}[subsection]{Situation}

\theoremstyle{remark}
\newtheorem{remark}[subsection]{Remark}
\newtheorem{remarks}[subsection]{Remarks}

\numberwithin{equation}{subsection}

% Macros
%
\def\lim{\mathop{\mathrm{lim}}\nolimits}
\def\colim{\mathop{\mathrm{colim}}\nolimits}
\def\Spec{\mathop{\mathrm{Spec}}}
\def\Hom{\mathop{\mathrm{Hom}}\nolimits}
\def\Ext{\mathop{\mathrm{Ext}}\nolimits}
\def\SheafHom{\mathop{\mathcal{H}\!\mathit{om}}\nolimits}
\def\SheafExt{\mathop{\mathcal{E}\!\mathit{xt}}\nolimits}
\def\Sch{\mathit{Sch}}
\def\Mor{\mathop{\mathrm{Mor}}\nolimits}
\def\Ob{\mathop{\mathrm{Ob}}\nolimits}
\def\Sh{\mathop{\mathit{Sh}}\nolimits}
\def\NL{\mathop{N\!L}\nolimits}
\def\CH{\mathop{\mathrm{CH}}\nolimits}
\def\proetale{{pro\text{-}\acute{e}tale}}
\def\etale{{\acute{e}tale}}
\def\QCoh{\mathit{QCoh}}
\def\Ker{\mathop{\mathrm{Ker}}}
\def\Im{\mathop{\mathrm{Im}}}
\def\Coker{\mathop{\mathrm{Coker}}}
\def\Coim{\mathop{\mathrm{Coim}}}

% Boxtimes
%
\DeclareMathSymbol{\boxtimes}{\mathbin}{AMSa}{"02}

%
% Macros for moduli stacks/spaces
%
\def\QCohstack{\mathcal{QC}\!\mathit{oh}}
\def\Cohstack{\mathcal{C}\!\mathit{oh}}
\def\Spacesstack{\mathcal{S}\!\mathit{paces}}
\def\Quotfunctor{\mathrm{Quot}}
\def\Hilbfunctor{\mathrm{Hilb}}
\def\Curvesstack{\mathcal{C}\!\mathit{urves}}
\def\Polarizedstack{\mathcal{P}\!\mathit{olarized}}
\def\Complexesstack{\mathcal{C}\!\mathit{omplexes}}
% \Pic is the operator that assigns to X its picard group, usage \Pic(X)
% \Picardstack_{X/B} denotes the Picard stack of X over B
% \Picardfunctor_{X/B} denotes the Picard functor of X over B
\def\Pic{\mathop{\mathrm{Pic}}\nolimits}
\def\Picardstack{\mathcal{P}\!\mathit{ic}}
\def\Picardfunctor{\mathrm{Pic}}
\def\Deformationcategory{\mathcal{D}\!\mathit{ef}}

\setlength{\emergencystretch}{3em}
\begin{document}
\title[Stacks Project, Tag 0208]{1668 --- Properness of finite-type morphisms to locally Noetherian schemes is detected by unique DVR lifts in generic-component function fields}
\maketitle
\noindent Copyright (C) 2005--2025 Johan de Jong. Stacks Project authors. Source: \href{https://stacks.math.columbia.edu/tag/0208}{Tag 0208}.

\noindent Editorial difficulty note (not author text). Difficulty H2: The strongest criterion tests only DVRs inside generic-component function fields, not arbitrary valuation rings. The author constructs a forbidden specialization on the projective closure of a Chow modification, produces a same-function-field DVR by Krull\textendash{}Akizuki, and transfers existence/uniqueness through the proper cover and diagonal. This geometric reduction is substantially beyond applying the usual valuative criterion..

\noindent Editorial provenance note (not author text). History: excerpt from revision \texttt{a0444\allowbreak 6e57e\allowbreak c1fbc\allowbreak 252a8\allowbreak 71afc\allowbreak ec775\allowbreak 2fb28\allowbreak 07b14}, limits.tex, lines 3926--4006. Reference presentation was adapted; original-author context and core support blocks were retained or added. The mathematical wording is unchanged. Licensed under GNU FDL 1.2 or later; no invariant sections or cover texts. See ../licenses/COPYING. Context blocks reproduce author definitions and supporting results; they are not additional counted problems.

\begin{situation}
\label{situation-descent}
Let $S = \lim_{i \in I} S_i$ be the limit of a directed system of schemes
with affine transition morphisms $f_{i'i} : S_{i'} \to S_i$
(Lemma \ref{lemma-directed-inverse-system-has-limit}).
We assume that $S_i$ is quasi-compact and quasi-separated for all $i \in I$.
We denote $f_i : S \to S_i$ the projection.
We also choose an element $0 \in I$.
\end{situation}

\begin{definition}
\label{properties-definition-ample}
\begin{reference}
\href{https://stacks.math.columbia.edu/bibliography}{Bibliography: EGA (II Definition 4.5.3)}
\end{reference}
Let $X$ be a scheme.
Let $\mathcal{L}$ be an invertible $\mathcal{O}_X$-module.
We say $\mathcal{L}$ is {\it ample} if
\begin{enumerate}
\item $X$ is quasi-compact, and
\item for every $x \in X$ there exists an $n \geq 1$
and $s \in \Gamma(X, \mathcal{L}^{\otimes n})$ such
that $x \in X_s$ and $X_s$ is affine.
\end{enumerate}
\end{definition}

\begin{definition}
\label{schemes-definition-reduced-induced-scheme}
Let $X$ be a scheme. Let $Z \subset X$ be a closed subset.
A {\it scheme structure on $Z$} is given by a closed subscheme $Z'$ of
$X$ whose underlying set is equal to $Z$. We often say
``let $(Z, \mathcal{O}_Z)$ be a scheme structure on $Z$'' to
indicate this. The {\it reduced induced scheme structure}
on $Z$ is the one constructed in Lemma \href{https://stacks.math.columbia.edu/tag/01J3}{lemma reduced closed subscheme (Tag 01J3)}.
The {\it reduction $X_{red}$ of $X$} is the reduced induced scheme
structure on $X$ itself.
\end{definition}

\begin{situation}
\label{situation-descent-property}
Let $S = \lim S_i$ be a limit of a directed system of schemes
with affine transition morphisms
(Lemma \ref{lemma-directed-inverse-system-has-limit}).
Let $0 \in I$ and let $f_0 : X_0 \to Y_0$ be a morphism of schemes over $S_0$.
Assume $S_0$, $X_0$, $Y_0$ are quasi-compact and quasi-separated.
Let $f_i : X_i \to Y_i$ be the base change of $f_0$ to $S_i$ and
let $f : X \to Y$ be the base change of $f_0$ to $S$.
\end{situation}

\begin{proposition}
\label{proposition-approximate}
Let $S$ be a quasi-compact and quasi-separated scheme.
There exist a directed set $I$
and an inverse system of schemes $(S_i, f_{ii'})$ over $I$
such that
\begin{enumerate}
\item the transition morphisms $f_{ii'}$ are affine
\item each $S_i$ is of finite type over $\mathbf{Z}$, and
\item $S = \lim_i S_i$.
\end{enumerate}
\end{proposition}

\begin{proposition}
\label{proposition-separated-closed-in-finite-presentation}
Let $f : X \to S$ be a morphism of schemes. Assume
\begin{enumerate}
\item $f$ is of finite type and separated, and
\item $S$ is quasi-compact and quasi-separated.
\end{enumerate}
Then there exists a separated morphism of finite presentation
$f' : X' \to S$ and a closed immersion $X \to X'$ of schemes over $S$.
\end{proposition}

\begin{definition}
\label{morphisms-definition-scheme-theoretic-image}
Let $f : X \to Y$ be a morphism of schemes. The {\it scheme theoretic image}
of $f$ is the smallest closed subscheme $Z \subset Y$ through which $f$
factors, see Lemma \href{https://stacks.math.columbia.edu/tag/01R6}{lemma scheme theoretic image (Tag 01R6)} above.
\end{definition}

\begin{definition}
\label{morphisms-definition-scheme-theoretically-dense}
Let $X$ be a scheme. Let $U \subset X$ be an open subscheme.
\begin{enumerate}
\item The scheme theoretic image of the morphism $U \to X$
is called the {\it scheme theoretic closure of $U$ in $X$}.
\item We say $U$ is {\it scheme theoretically dense in $X$}
if for every open $V \subset X$ the scheme theoretic closure
of $U \cap V$ in $V$ is equal to $V$.
\end{enumerate}
\end{definition}

\noindent
First we state the result concerning
separation. We will often use solid commutative diagrams of morphisms of
schemes having the following shape
\begin{equation}
\label{equation-valuative}
\vcenter{
\xymatrix{
\Spec(K) \ar[r] \ar[d] & X \ar[d] \\
\Spec(A) \ar[r] \ar@{-->}[ru] & S
}
}
\end{equation}
with $A$ a valuation ring and $K$ its field of fractions.

\begin{lemma}
\label{lemma-Noetherian-dvr-valuative-separation}
Let $S$ be a locally Noetherian scheme.
Let $f : X \to S$ be a morphism of schemes.
Assume $f$ is locally of finite type.
The following are equivalent:
\begin{enumerate}
\item The morphism $f$ is separated.
\item For any diagram (\ref{equation-valuative}) there is at most
one dotted arrow.
\item For all diagrams (\ref{equation-valuative}) with $A$ a discrete
valuation ring there is at most one dotted arrow.
\item For any irreducible component $X_0$ of $X$ with
generic point $\eta \in X_0$, for any discrete valuation ring
$A \subset K = \kappa(\eta)$ with fraction field $K$ and any
diagram (\ref{equation-valuative}) such that
the morphism $\Spec(K) \to X$ is the canonical one
(see Schemes, Section \href{https://stacks.math.columbia.edu/tag/01J5}{section points (Tag 01J5)})
there is at most one dotted arrow.
\end{enumerate}
\end{lemma}

\begin{proof}
Clearly (1) implies (2), (2) implies (3), and (3) implies (4).  It
remains to show (4) implies (1). Assume (4).
We begin by reducing to $S$ affine.  Being separated is a local
on the base (see
Schemes, Lemma \href{https://stacks.math.columbia.edu/tag/01KP}{lemma characterize separated (Tag 01KP)}).
Hence, if we can show that whenever
$X \to S$ has (4) that the restriction $X_\alpha \to S_\alpha$ has (4)
where $S_\alpha \subset S$ is an (affine) open subset and $X_\alpha :=
f^{-1}(S_\alpha)$, then we will be done.  The
generic points of the irreducible components of $X_\alpha$ will be the
generic points of irreducible components of $X$, since $X_\alpha$ is
open in $X$.  Therefore, any two distinct dotted arrows in the diagram
\begin{equation}
\label{equation-valuative-alpha}
\xymatrix{
\Spec(K) \ar[r] \ar[d] & X_\alpha \ar[d] \\
\Spec(A) \ar[r] \ar@{-->}[ru] & S_\alpha
}
\end{equation}
would then give two distinct arrows in diagram
(\ref{equation-valuative}) via the maps $X_\alpha \to X$ and
$S_\alpha \to S$, which is a contradiction.  Thus we have reduced
to the case $S$ is affine.  We remark that in the course of this
reduction, we prove that if $X \to S$ has (4) then the restriction $U
\to V$ has (4) for opens $U \subset X$ and $V \subset S$ with
$f(U) \subset V$.

\medskip\noindent
We next wish to reduce to the case $X \to S$ is finite type.  Assume
that we know (4) implies (1) when $X$ is finite type. Since
$S$ is Noetherian and $X$ is locally of finite type over $S$
we see $X$ is locally Noetherian as well (see Morphisms,
Lemma \href{https://stacks.math.columbia.edu/tag/01T6}{lemma finite type noetherian (Tag 01T6)}).
Thus, $X \to S$ is quasi-separated (see
Properties, Lemma \href{https://stacks.math.columbia.edu/tag/01OY}{lemma locally Noetherian quasi separated (Tag 01OY)}),
and therefore we may apply the valuative criterion to check whether $X$
is separated (see
Schemes, Lemma \href{https://stacks.math.columbia.edu/tag/01L0}{lemma valuative criterion separatedness (Tag 01L0)}).
Let $X = \bigcup_\alpha X_\alpha$ be an affine open
cover of $X$. Given any two dotted arrows, in a diagram
(\ref{equation-valuative}), the image of the closed points of
$\Spec A$ will
fall in two sets $X_\alpha$ and $X_\beta$.  Since $X_\alpha \cup
X_\beta$ is open, for topological reasons it must contain the image of
$\Spec(A)$ under both maps. Therefore, the two dotted arrows factor
through $X_\alpha \cup X_\beta \to X$, which is a scheme of finite type over
$S$. Since $X_\alpha \cup X_\beta$ is an open subset of $X$, by our
previous remark, $X_\alpha \cup X_\beta$ satisfies (4), so by
assumption, is separated.  This implies the two given dotted
arrows are the same. Therefore, we have reduced to $X \to S$ is finite type.

\medskip\noindent
Assume $X \to S$ of finite type and assume (4).
Since $X \to S$ is finite type, and $S$ is an affine Noetherian
scheme, $X$ is also Noetherian (see
Morphisms, Lemma \href{https://stacks.math.columbia.edu/tag/01T6}{lemma finite type noetherian (Tag 01T6)}).
Therefore, $X \to X \times_S X$ will
be a quasi-compact immersion of Noetherian schemes.  We proceed by
contradiction.  Assume that $X \to X \times_S X$ is not closed.  Then,
there is some $y \in X \times_S X$ in the closure of the image that is
not in the image. As $X$ is Noetherian it has finitely many irreducible
components. Therefore, $y$ is in the closure of the image of one of
the irreducible components $X_0 \subset X$.  Give $X_0$ the reduced
induced structure.  The composition $X_0 \to X \to X \times_S X$
factors through the closed subscheme $X_0 \times_S X_0 \subset X \times_S X$.
Denote the closure of $\Delta(X_0)$ in $X_0 \times_S X_0$
by $\bar X_0$ (again as a reduced closed subscheme). Thus $y \in \bar X_0$.
Since $X_0 \to X_0 \times_S X_0$ is an immersion, the image of $X_0$
will be open in $\bar X_0$. Hence $X_0$ and $\bar X_0$ are
birational. Since $\bar{X}_0$ is a closed subscheme of a
Noetherian scheme, it is Noetherian. Thus, the local ring
$\mathcal O_{{\bar X_0, y}}$ is a local Noetherian domain with fraction
field $K$ equal to the function field of $X_0$.  By the Krull-Akizuki
theorem (see Algebra, Lemma \ref{algebra-lemma-exists-dvr}), there exists a
discrete valuation ring $A$ dominating $\mathcal O_{{\bar X_0, y}}$
with fraction field $K$.  This allows to construct a diagram:
\begin{equation}
\label{equation-valuative-generic}
\xymatrix{
\Spec(K) \ar[r] \ar[d] & X_0 \ar[d]^{\Delta} \\
\Spec(A) \ar[r] \ar@{-->}[ur]& X_0 \times_S X_0 \\
}
\end{equation}
which sends $\Spec K$ to the generic point of $\Delta(X_0)$ and
the closed point of $A$ to $y \in X_0 \times_S X_0$ (use the material in
Schemes, Section \href{https://stacks.math.columbia.edu/tag/01J5}{section points (Tag 01J5)} to construct the arrows).
There cannot even exist
a set theoretic dotted arrow, since $y$ is not in the image of
$\Delta$ by our choice of $y$.  By categorical means, the existence of
the dotted arrow in the above diagram is equivalent to the uniqueness
of the dotted arrow in the following diagram:
\begin{equation}
\label{equation-valuative-nonexistent}
\xymatrix{
\Spec(K) \ar[r] \ar[d] & X_0 \ar[d]\\
\Spec(A) \ar[r] \ar@{-->}[ur] & S \\
}
\end{equation}
Therefore, we have non-uniqueness in this latter diagram by the
nonexistence in the first.  Therefore, $X_0$ does not satisfy
uniqueness for discrete valuation rings, and since $X_0$ is an
irreducible component of $X$, we have that $X \to S$ does not satisfy
(4).  Therefore, we have shown (4) implies (1).
\end{proof}


\begin{lemma}
\label{lemma-limited-base-change}
Let $S$ be a scheme.
Let $f : X \to S$ be a separated morphism of finite type.
The following are equivalent:
\begin{enumerate}
\item The morphism $f$ is proper.
\item For any morphism $S' \to S$ which is locally of finite type
the base change $X_{S'} \to S'$ is closed.
\item For every $n \geq 0$ the morphism
$\mathbf{A}^n \times X \to \mathbf{A}^n \times S$ is closed.
\end{enumerate}
\end{lemma}

\begin{proof}[First proof]
In view of the fact that a proper morphism is the same thing as
a separated, finite type, and universally closed morphism, this
lemma is a special case of Lemma \href{https://stacks.math.columbia.edu/tag/05JX}{lemma test universally closed (Tag 05JX)}.
\end{proof}

\begin{proof}[Second proof]
Clearly (1) implies (2), and (2) implies (3), so we just need to show (3)
implies (1).
First we reduce to the case when $S$ is affine.  Assume that (3) implies (1)
when the base is affine.  Now let $f: X \to S$ be a separated morphism of
finite type.  Being proper is local on the base
(see Morphisms, Lemma \href{https://stacks.math.columbia.edu/tag/01W2}{lemma proper local on the base (Tag 01W2)}), so if
$S = \bigcup_\alpha S_\alpha$ is an open affine cover, and if
we denote $X_\alpha := f^{-1}(S_\alpha)$, then it is
enough to show that $f|_{X_\alpha}: X_\alpha \to S_\alpha$ is proper
for all $\alpha$.  Since $S_\alpha$ is affine, if
the map $f|_{X_\alpha}$ satisfies (3), then it will satisfy (1)
by assumption, and will be proper.  To finish the reduction to the
case $S$ is affine, we must show that if $f: X \to S$ is separated of
finite type satisfying (3), then $f|_{X_\alpha} : X_\alpha \to S_\alpha$
is separated of finite type satisfying (3).  Separatedness and finite
type are clear.  To see (3), notice that
$\mathbf{A}^n \times X_\alpha$ is the open preimage of
$\mathbf{A}^n \times S_\alpha$ under the map $1 \times f$.  Fix a closed
set $Z \subset \mathbf A^n \times X_\alpha$.  Let $\bar Z$ denote the
closure of $Z$ in $\mathbf{A}^n \times X$.  Then for topological
reasons,
$$
1 \times f(\bar Z) \cap \mathbf{A}^n \times S_\alpha  = 1 \times f(Z).
$$
Hence $1 \times f(Z)$ is closed, and we have reduced the proof of
(3) $\Rightarrow$ (1) to the affine case.

\medskip\noindent
Assume $S$ affine, and $f : X \to S$ separated of finite type.
We can apply Chow's Lemma \ref{lemma-chow-finite-type}
to get $\pi : X' \to X$ proper surjective and $X' \to \mathbf{P}^n_S$
an immersion. If $X$ is proper over $S$, then $X' \to S$ is proper
(Morphisms, Lemma \href{https://stacks.math.columbia.edu/tag/01W3}{lemma composition proper (Tag 01W3)}). Since
$\mathbf{P}^n_S \to S$ is separated, we conclude that $X' \to
\mathbf{P}^n_S$ is proper
(Morphisms, Lemma \href{https://stacks.math.columbia.edu/tag/01W6}{lemma image proper scheme closed (Tag 01W6)})
and hence a closed immersion
(Schemes, Lemma \href{https://stacks.math.columbia.edu/tag/01IQ}{lemma immersion when closed (Tag 01IQ)}).
Conversely, assume $X' \to \mathbf{P}^n_S$ is a closed immersion.
Consider the diagram:
\begin{equation}
\label{equation-check-proper}
\xymatrix{
X' \ar[r] \ar@{->>}[d]_{\pi} &
\mathbf{P}^n_S \ar[d] \\
X \ar[r]^f & S
}
\end{equation}
All maps are a priori proper except for $X \to S$.
Hence we conclude that $X \to S$ is proper by
Morphisms, Lemma \href{https://stacks.math.columbia.edu/tag/03GN}{lemma image proper is proper (Tag 03GN)}.
Therefore, we have shown that $X \to S$ is proper if and only if
$X' \to \mathbf{P}^n_S$ is a closed immersion.

\medskip\noindent
Assume $S$ is affine and (3) holds, and let $n, X', \pi$ be as above.
Since being a closed morphism is local on the base, the map
$X \times \mathbf{P}^n \to S \times \mathbf{P}^n$ is closed since by (3)
$X \times \mathbf{A}^n \to S \times \mathbf{A}^n$ is closed and since
projective space is covered by copies of affine $n$-space, see
Constructions,
Lemma \href{https://stacks.math.columbia.edu/tag/01NG}{lemma standard covering projective space (Tag 01NG)}.
By Morphisms, Lemma \href{https://stacks.math.columbia.edu/tag/01W4}{lemma base change proper (Tag 01W4)}
the morphism
$$
X' \times_S \mathbf{P}^n_S
\to
X \times_S \mathbf{P}^n_S =
X \times \mathbf{P}^n
$$
is proper. Since $\mathbf{P}^n$ is separated, the projection
$$
X' \times_S \mathbf{P}^n_S = \mathbf{P}^n_{X'} \to X'
$$
will be separated as it is just a base change of a separated
morphism. Therefore, the map $X' \to X' \times_S \mathbf{P}^n_S$ is proper,
since it is a section to a separated map (see
Schemes, Lemma \href{https://stacks.math.columbia.edu/tag/01KT}{lemma section immersion (Tag 01KT)}).
Composing these morphisms
$$
X' \to X' \times_S \mathbf{P}^n_S \to X \times_S \mathbf{P}^n_S
= X \times \mathbf{P}^n \to S \times \mathbf{P}^n = \mathbf{P}^n_S
$$
we find that the immersion $X' \to \mathbf{P}^n_S$ is closed, and
hence a closed immersion.
\end{proof}


\begin{lemma}
\label{lemma-chow-finite-type}
Let $S$ be a quasi-compact and quasi-separated scheme.
Let $f : X \to S$ be a separated morphism of finite type.
Then there exists an $n \geq 0$ and a diagram
$$
\xymatrix{
X \ar[rd] & X' \ar[d] \ar[l]^\pi \ar[r] & \mathbf{P}^n_S \ar[dl] \\
& S &
}
$$
where $X' \to \mathbf{P}^n_S$ is an immersion, and
$\pi : X' \to X$ is proper and surjective.
\end{lemma}

\begin{proof}
By Proposition \ref{proposition-separated-closed-in-finite-presentation}
we can find a closed immersion $X \to Y$ where $Y$ is separated
and of finite presentation over $S$. Clearly, if we prove the assertion
for $Y$, then the result follows for $X$. Hence we may assume that
$X$ is of finite presentation over $S$.

\medskip\noindent
Write $S = \lim_i S_i$ as a directed limit of Noetherian schemes, see
Proposition \ref{proposition-approximate}. By
Lemma \ref{lemma-descend-finite-presentation} we can
find an index $i \in I$ and a scheme $X_i \to S_i$ of finite presentation
so that $X = S \times_{S_i} X_i$.
By Lemma \ref{lemma-descend-separated-finite-presentation}
we may assume that $X_i \to S_i$ is separated.
Clearly, if we prove the assertion for
$X_i$ over $S_i$, then the assertion holds for $X$. The case
$X_i \to S_i$ is treated by
Cohomology of Schemes, Lemma \ref{coherent-lemma-chow-Noetherian}.
\end{proof}


\begin{lemma}
\label{coherent-lemma-chow-Noetherian}
\begin{reference}
\href{https://stacks.math.columbia.edu/bibliography}{Bibliography: EGA (II Theorem 5.6.1(a))}
\end{reference}
Let $S$ be a Noetherian scheme.
Let $f : X \to S$ be a separated morphism of finite type.
Then there exist an $n \geq 0$ and a diagram
$$
\xymatrix{
X \ar[rd] & X' \ar[d] \ar[l]^\pi \ar[r] & \mathbf{P}^n_S \ar[dl] \\
& S &
}
$$
where $X' \to \mathbf{P}^n_S$ is an immersion, and
$\pi : X' \to X$ is proper and surjective. Moreover, we may
arrange it such that there exists a dense open subscheme
$U \subset X$ such that $\pi^{-1}(U) \to U$ is an isomorphism.
\end{lemma}

\begin{proof}
All of the schemes we will encounter during the rest of the proof
are going to be of finite type over the Noetherian scheme $S$ and
hence Noetherian
(see Morphisms, Lemma \href{https://stacks.math.columbia.edu/tag/01T6}{lemma finite type noetherian (Tag 01T6)}).
All morphisms between them will automatically be quasi-compact, locally of
finite type and quasi-separated, see
Morphisms, Lemma \href{https://stacks.math.columbia.edu/tag/01T8}{lemma permanence finite type (Tag 01T8)} and
Properties,
Lemmas \href{https://stacks.math.columbia.edu/tag/01OY}{lemma locally Noetherian quasi separated (Tag 01OY)} and
\href{https://stacks.math.columbia.edu/tag/01P0}{lemma morphism Noetherian schemes quasi compact (Tag 01P0)}.

\medskip\noindent
The scheme $X$ has only finitely many irreducible components
(Properties, Lemma \href{https://stacks.math.columbia.edu/tag/0BA8}{lemma Noetherian irreducible components (Tag 0BA8)}).
Say $X = X_1 \cup \ldots \cup X_r$ is the decomposition
of $X$ into irreducible components.
Let $\eta_i \in X_i$ be the generic point.
For every point $x \in X$ there exists an affine open
$U_x \subset X$ which contains $x$ and each of the generic
points $\eta_i$. See
Properties, Lemma \href{https://stacks.math.columbia.edu/tag/01ZX}{lemma point and maximal points affine (Tag 01ZX)}.
Since $X$ is quasi-compact, we can find a finite affine open
covering $X = U_1 \cup \ldots \cup U_m$ such that
each $U_i$ contains $\eta_1, \ldots, \eta_r$.
In particular we conclude that the open
$U = U_1 \cap \ldots \cap U_m \subset X$ is
a dense open. This and the fact that the $U_i$ are affine opens
covering $X$ are all that we will use below.

\medskip\noindent
Let $X^* \subset X$ be the scheme theoretic closure of $U \to X$, see
Morphisms, Definition \ref{morphisms-definition-scheme-theoretic-image}.
Let $U_i^* = X^* \cap U_i$. Note that $U_i^*$ is a closed subscheme
of $U_i$. Hence $U_i^*$ is affine. Since $U$ is dense in $X$ the
morphism $X^* \to X$ is a surjective closed immersion. It is an
isomorphism over $U$. Hence we may replace $X$ by $X^*$ and
$U_i$ by $U_i^*$ and assume that $U$ is scheme theoretically dense
in $X$, see
Morphisms, Definition \ref{morphisms-definition-scheme-theoretically-dense}.

\medskip\noindent
By Morphisms, Lemma \href{https://stacks.math.columbia.edu/tag/01VS}{lemma quasi projective finite type over S (Tag 01VS)}
we can find an immersion $j_i : U_i \to \mathbf{P}_S^{n_i}$
for each $i$. By
Morphisms, Lemma \href{https://stacks.math.columbia.edu/tag/01RG}{lemma quasi compact immersion (Tag 01RG)} we can find
closed subschemes $Z_i \subset \mathbf{P}_S^{n_i}$
such that $j_i : U_i \to Z_i$ is a scheme theoretically
dense open immersion. Note that $Z_i \to S$ is proper, see
Morphisms, Lemma \href{https://stacks.math.columbia.edu/tag/01WC}{lemma locally projective proper (Tag 01WC)}.
Consider the morphism
$$
j = (j_1|_U, \ldots, j_m|_U) : U \longrightarrow
\mathbf{P}_S^{n_1} \times_S \ldots \times_S \mathbf{P}_S^{n_m}.
$$
By the lemma cited above we can find a closed subscheme
$Z$ of $\mathbf{P}_S^{n_1} \times_S \ldots \times_S \mathbf{P}_S^{n_m}$
such that $j : U \to Z$ is an open immersion and such that $U$
is scheme theoretically dense in $Z$. The morphism $Z \to S$
is proper. Consider the $i$th projection
$$
\text{pr}_i|_Z : Z \longrightarrow \mathbf{P}^{n_i}_S.
$$
This morphism factors through $Z_i$ (see Morphisms,
Lemma \href{https://stacks.math.columbia.edu/tag/01R9}{lemma factor factor (Tag 01R9)}). Denote $p_i : Z \to Z_i$
the induced morphism. This is a proper morphism, see
Morphisms, Lemma \href{https://stacks.math.columbia.edu/tag/01W6}{lemma image proper scheme closed (Tag 01W6)}
for example. At this point we have that
$U \subset U_i \subset Z_i$ are scheme theoretically
dense open immersions. Moreover, we can think of $Z$ as the
scheme theoretic image of the ``diagonal'' morphism
$U \to Z_1 \times_S \ldots \times_S Z_m$.

\medskip\noindent
Set $V_i = p_i^{-1}(U_i)$. Note that $p_i|_{V_i} : V_i \to U_i$ is proper.
Set $X' = V_1 \cup \ldots \cup V_m$. By construction $X'$ has an immersion
into the scheme
$\mathbf{P}^{n_1}_S \times_S \ldots \times_S \mathbf{P}^{n_m}_S$.
Thus by the Segre embedding (see
Constructions, Lemma \href{https://stacks.math.columbia.edu/tag/01WD}{lemma segre embedding (Tag 01WD)})
we see that $X'$ has
an immersion into a projective space over $S$.

\medskip\noindent
We claim that the morphisms $p_i|_{V_i}: V_i \to U_i$ glue to a morphism
$X' \to X$. Namely, it is clear that $p_i|_U$ is the identity map
from $U$ to $U$. Since $U \subset X'$ is scheme theoretically
dense by construction, it is also scheme theoretically dense
in the open subscheme $V_i \cap V_j$. Thus we see that
$p_i|_{V_i \cap V_j} = p_j|_{V_i \cap V_j}$ as morphisms into the
separated $S$-scheme $X$, see
Morphisms, Lemma \href{https://stacks.math.columbia.edu/tag/01RH}{lemma equality of morphisms (Tag 01RH)}.
We denote the resulting morphism $\pi : X' \to X$.

\medskip\noindent
We claim that $\pi^{-1}(U_i) = V_i$.
Since $\pi|_{V_i} = p_i|_{V_i}$ it follows that
$V_i \subset \pi^{-1}(U_i)$. Consider the diagram
$$
\xymatrix{
V_i \ar[r] \ar[rd]_{p_i|_{V_i}} & \pi^{-1}(U_i) \ar[d] \\
& U_i
}
$$
Since $V_i \to U_i$ is proper we see that the image of
the horizontal arrow is closed, see
Morphisms, Lemma \href{https://stacks.math.columbia.edu/tag/01W6}{lemma image proper scheme closed (Tag 01W6)}.
Since $V_i \subset \pi^{-1}(U_i)$ is scheme
theoretically dense (as it contains $U$)
we conclude that $V_i = \pi^{-1}(U_i)$ as claimed.

\medskip\noindent
This shows that $\pi^{-1}(U_i) \to U_i$ is identified with the proper
morphism $p_i|_{V_i} : V_i \to U_i$. Hence we see that $X$ has a
finite affine covering $X = \bigcup U_i$ such that the restriction
of $\pi$ is proper on each member of the covering.
Thus by Morphisms, Lemma \href{https://stacks.math.columbia.edu/tag/01W2}{lemma proper local on the base (Tag 01W2)}
we see that $\pi$ is proper.

\medskip\noindent
Finally we have to show that $\pi^{-1}(U) = U$. To see this we argue in the
same way as above using the diagram
$$
\xymatrix{
U \ar[r] \ar[rd] & \pi^{-1}(U) \ar[d] \\
& U
}
$$
and using that $\text{id}_U : U \to U$ is proper and that
$U$ is scheme theoretically dense in $\pi^{-1}(U)$.
\end{proof}


\begin{lemma}
\label{algebra-lemma-exists-dvr}
Let $R$ be a Noetherian local domain with fraction field $K$.
Assume that $R$ is not a field.
Let $L/K$ be a finitely generated field extension.
Then there exists discrete valuation ring $A$ with fraction field
$L$ which dominates $R$.
\end{lemma}

\begin{proof}
If $L$ is not finite over $K$ choose a transcendence basis
$x_1, \ldots, x_r$ of $L$ over $K$ and replace $R$ by
$R[x_1, \ldots, x_r]$ localized at the maximal ideal
generated by $\mathfrak m_R$ and $x_1, \ldots, x_r$.
Thus we may assume $K \subset L$ finite.

\medskip\noindent
By Lemma \ref{algebra-lemma-dominate-by-dimension-1} we may assume $\dim(R) = 1$.

\medskip\noindent
Let $A \subset L$ be the integral closure of $R$ in $L$.
By Lemma \ref{algebra-lemma-krull-akizuki} this is Noetherian.
By Lemma \href{https://stacks.math.columbia.edu/tag/00GQ}{lemma integral overring surjective (Tag 00GQ)} there
is a prime ideal $\mathfrak q \subset A$ lying
over the maximal ideal of $R$.
By Lemma \href{https://stacks.math.columbia.edu/tag/00PD}{lemma characterize dvr (Tag 00PD)} the ring $A_{\mathfrak q}$ is a discrete
valuation ring dominating $R$ as desired.
\end{proof}


\begin{lemma}
\label{algebra-lemma-dominate-by-dimension-1}
Let $R$ be a local Noetherian domain with fraction field $K$.
Assume $R$ is not a field.
Then there exist $R \subset R' \subset K$ with
\begin{enumerate}
\item $R'$ local Noetherian of dimension $1$,
\item $R \to R'$ a local ring map, i.e., $R'$ dominates $R$, and
\item $R \to R'$ essentially of finite type.
\end{enumerate}
\end{lemma}

\begin{proof}
Choose any valuation ring $A \subset K$ dominating $R$ (which
exist by Lemma \href{https://stacks.math.columbia.edu/tag/00IA}{lemma dominate (Tag 00IA)}).
Denote $v$ the corresponding valuation.
Let $x_1, \ldots, x_r$ be a minimal set of generators of the
maximal ideal $\mathfrak m$ of $R$. We may and do assume that
$v(x_r) = \min\{v(x_1), \ldots, v(x_r)\}$. Consider the ring
$$
S = R[x_1/x_r, x_2/x_r, \ldots, x_{r - 1}/x_r] \subset K.
$$
Note that $\mathfrak mS = x_rS$ is a principal ideal.
Note that $S \subset A$ and that $v(x_r) > 0$, hence we see
that $x_rS \not = S$. Choose a minimal prime $\mathfrak q$
over $x_rS$. Then $\text{height}(\mathfrak q) = 1$ by
Lemma \href{https://stacks.math.columbia.edu/tag/00KV}{lemma minimal over 1 (Tag 00KV)}
and $\mathfrak q$ lies over $\mathfrak m$. Hence we
see that $R' = S_{\mathfrak q}$ is a solution.
\end{proof}


\begin{lemma}[Krull-Akizuki]
\label{algebra-lemma-krull-akizuki}
Let $R$ be a domain with fraction field $K$.
Let $L/K$ be a finite extension of fields.
Assume $R$ is Noetherian and $\dim(R) = 1$.
In this case any ring $A$ with $R \subset A \subset L$ is
Noetherian.
\end{lemma}

\begin{proof}
Let $I \subset A$ be a nonzero ideal. By
Lemma \href{https://stacks.math.columbia.edu/tag/0H7L}{lemma intersection not zero (Tag 0H7L)}
we can find a nonzero element $x \in I \cap R$.
Then we get $I/xA \subset A/xA$. By Lemma \href{https://stacks.math.columbia.edu/tag/00PF}{lemma finite length global (Tag 00PF)}
the $R$-module $A/xA$ has finite length as an $R$-module. Hence
$I/xA$ has finite length as an $R$-module. Hence $I$ is finitely
generated as an ideal in $A$.
\end{proof}


\begin{lemma}
\label{lemma-approximate}
Let $S$ be a quasi-compact and quasi-separated scheme. Let $V \subset S$
be a quasi-compact open. Let $I$ be a directed set
and let $(V_i, f_{ii'})$ be an inverse system of schemes over $I$
with affine transition maps, with each $V_i$ of finite type over $\mathbf{Z}$,
and with $V = \lim V_i$. Then there exist
\begin{enumerate}
\item a directed set $J$,
\item an inverse system of schemes $(S_j, g_{jj'})$ over $J$,
\item an order preserving map $\alpha : J \to I$,
\item open subschemes $V'_j \subset S_j$, and
\item isomorphisms $V'_j \to V_{\alpha(j)}$
\end{enumerate}
such that
\begin{enumerate}
\item the transition morphisms $g_{jj'} : S_j \to S_{j'}$ are affine,
\item each $S_j$ is of finite type over $\mathbf{Z}$,
\item $g_{jj'}^{-1}(V'_{j'}) = V'_j$,
\item $S = \lim S_j$ and $V = \lim V'_j$, and
\item the diagrams
$$
\vcenter{
\xymatrix{
V \ar[d] \ar[rd] \\
V'_j \ar[r] & V_{\alpha(j)}
}
}
\quad\text{and}\quad
\vcenter{
\xymatrix{
V'_j \ar[r] \ar[d] & V_{\alpha(j)} \ar[d] \\
V'_{j'} \ar[r] & V_{\alpha(j')}
}
}
$$
are commutative.
\end{enumerate}
\end{lemma}

\begin{proof}
Set $Z = S \setminus V$. Choose affine opens $U_1, \ldots, U_m \subset S$
such that $Z \subset \bigcup_{l = 1, \ldots, m} U_l$. Consider the opens
$$
V \subset V \cup U_1 \subset V \cup U_1 \cup U_2 \subset
\ldots \subset V \cup \bigcup\nolimits_{l = 1, \ldots, m} U_l = S
$$
If we can prove the lemma successively for each of the cases
$$
V \cup U_1 \cup \ldots \cup U_l
\subset
V \cup U_1 \cup \ldots \cup U_{l + 1}
$$
then the lemma will follow for $V \subset S$. In each case we are adding
one affine open. Thus we may assume
\begin{enumerate}
\item $S = U \cup V$,
\item $U$ affine open in $S$,
\item $V$ quasi-compact open in $S$, and
\item $V = \lim_i V_i$ with $(V_i, f_{ii'})$
an inverse system over a directed set $I$, each $f_{ii'}$
affine and each $V_i$ of finite type over $\mathbf{Z}$.
\end{enumerate}
Denote $f_i : V \to V_i$ the projections.
Set $W = U \cap V$. As $S$ is quasi-separated, this is a quasi-compact open
of $V$. By Lemma \ref{lemma-descend-opens}
(and after shrinking $I$) we may assume that there exist
opens $W_i \subset V_i$ such that $f_{ii'}^{-1}(W_{i'}) = W_i$
and such that $f_i^{-1}(W_i) = W$. Since $W$ is a quasi-compact open
of $U$ it is quasi-affine. Hence we may assume (after shrinking $I$ again)
that $W_i$ is quasi-affine for all $i$, see
Lemma \ref{lemma-limit-quasi-affine}.

\medskip\noindent
Write $U = \Spec(B)$. Set $R = \Gamma(W, \mathcal{O}_W)$,
and $R_i = \Gamma(W_i, \mathcal{O}_{W_i})$.
By Lemma \ref{lemma-descend-section} we have $R = \colim_i R_i$.
Now we have the maps of rings
$$
\xymatrix{
B \ar[r]_s & R \\
& R_i \ar[u]_{t_i}
}
$$
We set $B_i = \{(b, r) \in B \times R_i \mid s(b) = t_i(r)\}$ so that we
have a cartesian diagram
$$
\xymatrix{
B \ar[r]_s & R \\
B_i \ar[u] \ar[r] & R_i \ar[u]_{t_i}
}
$$
for each $i$. The transition maps $R_i \to R_{i'}$ induce maps
$B_i \to B_{i'}$. It is clear that $B = \colim_i B_i$.
In the next paragraph we show that for all sufficiently large $i$
the composition $W_i \to \Spec(R_i) \to \Spec(B_i)$ is an open immersion.

\medskip\noindent
As $W$ is a quasi-compact open of $U = \Spec(B)$
we can find a finitely many elements $g_l \in B$, $l = 1, \ldots, m$
such that $D(g_l) \subset W$ and such that
$W = \bigcup_{l = 1, \ldots, m} D(g_l)$.
Note that this implies $D(g_l) = W_{s(g_l)}$ as open subsets of $U$,
where $W_{s(g_l)}$ denotes the largest open subset of $W$ on which
$s(g_l)$ is invertible. Hence
$$
B_{g_l} =
\Gamma(D(g_l), \mathcal{O}_U) =
\Gamma(W_{s(g_l)}, \mathcal{O}_W) = R_{s(g_l)},
$$
where the last equality is
Properties, Lemma \href{https://stacks.math.columbia.edu/tag/01P7}{lemma invert f sections (Tag 01P7)}.
Since $W_{s(g_l)}$ is affine this also
implies that $D(s(g_l)) = W_{s(g_l)}$ as open subsets of $\Spec(R)$.
Since $R = \colim_i R_i$ we can (after shrinking $I$)
assume there exist $g_{l, i} \in R_i$ for all $i \in I$ such that
$s(g_l) = t_i(g_{l, i})$. Of course we choose the $g_{l, i}$
such that $g_{l, i}$ maps to $g_{l, i'}$ under the transition maps
$R_i \to R_{i'}$. Then, by Lemma \ref{lemma-descend-opens} we can
(after shrinking $I$ again)
assume the corresponding opens $D(g_{l, i}) \subset \Spec(R_i)$
are contained in $W_i$ for $l = 1, \ldots, m$ and cover $W_i$.
We conclude that the morphism $W_i \to \Spec(R_i) \to \Spec(B_i)$
is an open immersion, see Lemma \ref{lemma-diagram}.

\medskip\noindent
By Lemma \ref{lemma-quasi-affine-finite-type-over-Z}
we can write $B_i$ as a directed colimit of subalgebras
$A_{i, p} \subset B_i$, $p \in P_i$ each
of finite type over $\mathbf{Z}$ and such that $W_i$ is
identified with an open subscheme of $\Spec(A_{i, p})$.
Let $S_{i, p}$ be the scheme obtained by glueing
$V_i$ and $\Spec(A_{i, p})$ along the open $W_i$, see
Schemes, Section \href{https://stacks.math.columbia.edu/tag/01JA}{section glueing schemes (Tag 01JA)}.
Here is the resulting commutative diagram of schemes:
$$
\xymatrix{
& & V \ar[lld] \ar[d] & W \ar[l] \ar[lld] \ar[d] \\
V_i \ar[d] & W_i \ar[l] \ar[d] & S \ar[lld] & U \ar[lld] \ar[l] \\
S_{i, p} & \Spec(A_{i, p}) \ar[l]
}
$$
The morphism $S \to S_{i, p}$ arises because the upper right
square is a pushout in the category of schemes.
Note that $S_{i, p}$ is of finite type over $\mathbf{Z}$ since
it has a finite affine open covering whose members are
spectra of finite type $\mathbf{Z}$-algebras.
We define a preorder on $J = \coprod_{i \in I} P_i$
by the rule $(i', p') \geq (i, p)$ if and only if
$i' \geq i$ and the map $B_i \to B_{i'}$ maps $A_{i, p}$ into
$A_{i', p'}$. This is exactly the condition needed to
define a morphism $S_{i', p'} \to S_{i, p}$: namely make a commutative
diagram as above using the transition morphisms $V_{i'} \to V_i$
and $W_{i'} \to W_i$ and
the morphism $\Spec(A_{i', p'}) \to \Spec(A_{i, p})$ induced
by the ring map $A_{i, p} \to A_{i', p'}$. The relevant commutativities
have been built into the constructions.
We claim that $S$ is the directed limit of the schemes $S_{i, p}$.
Since by construction the schemes $V_i$ have limit $V$ this boils
down to the fact that $B$ is the limit of the rings $A_{i, p}$
which is true by construction. The map $\alpha : J \to I$ is given
by the rule $j = (i, p) \mapsto i$. The open subscheme $V'_j$ is
just the image of $V_i \to S_{i, p}$ above. The commutativity of
the diagrams in (5) is clear from the construction.
This finishes the proof of the lemma.
\end{proof}


\begin{lemma}
\label{lemma-descend-finite-presentation}
Let $I$ be a directed set.
Let $(S_i, f_{ii'})$ be an inverse system of schemes over $I$.
Assume
\begin{enumerate}
\item the morphisms $f_{ii'} : S_i \to S_{i'}$ are affine,
\item the schemes $S_i$ are quasi-compact and quasi-separated.
\end{enumerate}
Let $S = \lim_i S_i$. Then we have the following:
\begin{enumerate}
\item For any morphism of finite presentation $X \to S$
there exists an index $i \in I$ and a morphism of finite
presentation $X_i \to S_i$ such that $X \cong X_{i, S}$ as
schemes over $S$.
\item Given an index $i \in I$, schemes
$X_i$, $Y_i$ of finite presentation over $S_i$, and a morphism
$\varphi : X_{i, S} \to Y_{i, S}$ over $S$, there exists an index
$i' \geq i$ and a morphism
$\varphi_{i'} : X_{i, S_{i'}} \to Y_{i, S_{i'}}$
whose base change to $S$ is $\varphi$.
\item Given an index $i \in I$, schemes $X_i$, $Y_i$ of finite presentation
over $S_i$ and a pair of morphisms $\varphi_i, \psi_i : X_i \to Y_i$
whose base changes $\varphi_{i, S} = \psi_{i, S}$ are equal,
there exists an index $i' \geq i$ such that
$\varphi_{i, S_{i'}} = \psi_{i, S_{i'}}$.
\end{enumerate}
In other words, the category of schemes of finite presentation over
$S$ is the colimit over $I$ of the categories of schemes of finite
presentation over $S_i$.
\end{lemma}

\begin{proof}
In case each of the schemes $S_i$ is affine, and we consider
only affine schemes of finite presentation over $S_i$, resp.\ $S$
this lemma is equivalent to
Algebra, Lemma \ref{algebra-lemma-colimit-category-fp-algebras}.
We claim that the affine case implies the lemma in general.

\medskip\noindent
Let us prove (3). Suppose given an index $i \in I$, schemes
$X_i$, $Y_i$ of finite presentation over $S_i$ and a pair of morphisms
$\varphi_i, \psi_i : X_i \to Y_i$. Assume that the base changes are
equal: $\varphi_{i, S} = \psi_{i, S}$. We will use the notation
$X_{i'} = X_{i, S_{i'}}$ and $Y_{i'} = Y_{i, S_{i'}}$ for
$i' \geq i$. We also set $X = X_{i, S}$ and $Y = Y_{i, S}$.
Note that according to Lemma \ref{lemma-scheme-over-limit} we have
$X = \lim_{i' \geq i} X_{i'}$ and similarly for $Y$.
Additionally we denote $\varphi_{i'}$ and $\psi_{i'}$
(resp.\ $\varphi$ and $\psi$)
the base change of $\varphi_i$ and $\psi_i$ to $S_{i'}$
(resp.\ $S$). So our assumption means that $\varphi = \psi$.
Since $Y_i$ and $X_i$ are of finite presentation
over $S_i$, and since $S_i$ is quasi-compact and quasi-separated, also
$X_i$ and $Y_i$ are quasi-compact and quasi-separated
(see Morphisms,
Lemma \href{https://stacks.math.columbia.edu/tag/01TY}{lemma finite presentation quasi compact quasi separated (Tag 01TY)}).
Hence we may choose a finite affine open covering
$Y_i = \bigcup V_{j, i}$ such that each $V_{j, i}$ maps into
an affine open of $S_i$. As above, denote $V_{j, i'}$ the inverse
image of $V_{j, i}$ in $Y_{i'}$ and $V_j$ the inverse image in $Y$.
The immersions $V_{j, i'} \to Y_{i'}$ are quasi-compact, and the inverse images
$U_{j, i'} = \varphi_i^{-1}(V_{j, i'})$ and
$U_{j, i'}' = \psi_i^{-1}(V_{j, i'})$
are quasi-compact opens of $X_{i'}$. By assumption the inverse images of
$V_j$ under $\varphi$ and $\psi$ in $X$ are equal.
Hence by Lemma \ref{lemma-descend-opens}
there exists an index $i' \geq i$ such that
of $U_{j, i'} = U_{j, i'}'$ in $X_{i'}$.
Choose an finite affine open covering
$U_{j, i'} = U_{j, i'}' = \bigcup W_{j, k, i'}$
which induce coverings $U_{j, i''} = U_{j, i''}' = \bigcup W_{j, k, i''}$
for all $i'' \geq i'$.
By the affine case there exists
an index $i''$ such that
$\varphi_{i''}|_{W_{j, k, i''}} = \psi_{i''}|_{W_{j, k, i''}}$
for all $j, k$. Then $i''$ is an index such that
$\varphi_{i''} = \psi_{i''}$ and (3) is proved.

\medskip\noindent
Let us prove (2). Suppose given an index $i \in I$, schemes
$X_i$, $Y_i$ of finite presentation over $S_i$ and a morphism
$\varphi : X_{i, S} \to Y_{i, S}$. We will use the notation
$X_{i'} = X_{i, S_{i'}}$ and $Y_{i'} = Y_{i, S_{i'}}$ for
$i' \geq i$. We also set $X = X_{i, S}$ and $Y = Y_{i, S}$.
Note that according to Lemma \ref{lemma-scheme-over-limit} we have
$X = \lim_{i' \geq i} X_{i'}$ and similarly for $Y$.
Since $Y_i$ and $X_i$ are of finite presentation
over $S_i$, and since $S_i$ is quasi-compact and quasi-separated, also
$X_i$ and $Y_i$ are quasi-compact and quasi-separated
(see Morphisms,
Lemma \href{https://stacks.math.columbia.edu/tag/01TY}{lemma finite presentation quasi compact quasi separated (Tag 01TY)}).
Hence we may choose a finite affine open covering
$Y_i = \bigcup V_{j, i}$ such that each $V_{j, i}$ maps into
an affine open of $S_i$. As above, denote $V_{j, i'}$ the inverse
image of $V_{j, i}$ in $Y_{i'}$ and $V_j$ the inverse image in $Y$.
The immersions $V_j \to Y$ are quasi-compact, and the inverse images
$U_j = \varphi^{-1}(V_j)$ are quasi-compact opens of $X$.
Hence by Lemma \ref{lemma-descend-opens} there exists an index
$i' \geq i$ and quasi-compact opens $U_{j, i'}$ of $X_{i'}$
whose inverse image in $X$ is $U_j$. Choose an finite affine open covering
$U_{j, i'} = \bigcup W_{j, k, i'}$ which induce affine open
coverings $U_{j, i''} = \bigcup W_{j, k, i''}$
for all $i'' \geq i'$ and an affine open covering
$U_j = \bigcup W_{j, k}$. By the affine case there exists
an index $i''$ and morphisms
$\varphi_{j, k, i''} : W_{j, k, i''} \to V_{j, i''}$
such that
$\varphi|_{W_{j, k}} = \varphi_{j, k, i'', S}$ for all $j, k$.
By part (3) proved above, there is a further index $i''' \geq i''$
such that
$$
\varphi_{j_1, k_1, i'', S_{i'''}}|_{W_{j_1, k_1, i'''} \cap W_{j_2, k_2, i'''}}
=
\varphi_{j_2, k_2, i'', S_{i'''}}|_{W_{j_1, k_1, i'''} \cap W_{j_2, k_2, i'''}}
$$
for all $j_1, j_2, k_1, k_2$. Then $i'''$ is an index such that
there exists a morphism $\varphi_{i'''} : X_{i'''} \to Y_{i'''}$
whose base change to $S$ gives $\varphi$. Hence (2) holds.

\medskip\noindent
Let us prove (1). Suppose given a scheme $X$ of finite presentation
over $S$. Since $X$ is of finite presentation
over $S$, and since $S$ is quasi-compact and quasi-separated, also
$X$ is quasi-compact and quasi-separated
(see Morphisms,
Lemma \href{https://stacks.math.columbia.edu/tag/01TY}{lemma finite presentation quasi compact quasi separated (Tag 01TY)}).
Choose a finite affine open covering $X = \bigcup U_j$
such that each $U_j$ maps into an affine open $V_j \subset S$.
Denote $U_{j_1j_2} = U_{j_1} \cap U_{j_2}$ and
$U_{j_1j_2j_3} = U_{j_1} \cap U_{j_2} \cap U_{j_3}$.
By Lemmas \ref{lemma-descend-opens} and \ref{lemma-limit-affine}
we can find an index $i_1$ and affine opens $V_{j, i_1} \subset S_{i_1}$
such that each $V_j$ is the inverse of this in $S$.
Let $V_{j, i}$ be the inverse image of $V_{j, i_1}$ in $S_i$ for
$i \geq i_1$. By the affine case we may find an index $i_2 \geq i_1$ and
affine schemes $U_{j, i_2} \to V_{j, i_2}$ such
that $U_j = S \times_{S_{i_2}} U_{j, i_2}$ is the base change.
Denote $U_{j, i} = S_i \times_{S_{i_2}} U_{j, i_2}$ for $i \geq i_2$.
By Lemma \ref{lemma-descend-opens} there exists an index
$i_3 \geq i_2$ and open subschemes
$W_{j_1, j_2, i_3} \subset U_{j_1, i_3}$
whose base change to $S$ is equal to $U_{j_1j_2}$.
Denote $W_{j_1, j_2, i} = S_i \times_{S_{i_3}} W_{j_1, j_2, i_3}$
for $i \geq i_3$. By part (2) shown above there exists an index
$i_4 \geq i_3$ and morphisms
$\varphi_{j_1, j_2, i_4} : W_{j_1, j_2, i_4} \to W_{j_2, j_1, i_4}$
whose base change to $S$ gives the identity morphism
$U_{j_1j_2} = U_{j_2j_1}$ for all $j_1, j_2$.
For all $i \geq i_4$ denote
$\varphi_{j_1, j_2, i} = \text{id}_S \times \varphi_{j_1, j_2, i_4}$
the base change. We claim that for some $i_5 \geq i_4$ the system
$((U_{j, i_5})_j, (W_{j_1, j_2, i_5})_{j_1, j_2},
(\varphi_{j_1, j_2, i_5})_{j_1, j_2})$ forms a glueing datum
as in Schemes, Section \href{https://stacks.math.columbia.edu/tag/01JA}{section glueing schemes (Tag 01JA)}.
In order to see this we have to verify that for $i$ large enough
we have
$$
\varphi_{j_1, j_2, i}^{-1}(W_{j_2, j_1, i} \cap W_{j_2, j_3, i})
=
W_{j_1, j_2, i} \cap W_{j_1, j_3, i}
$$
and that for large enough $i$ the cocycle condition holds.
The first condition follows from Lemma \ref{lemma-descend-opens}
and the fact that $U_{j_2j_1j_3} = U_{j_1j_2j_3}$.
The second from part (3) of the lemma proved above and the fact
that the cocycle condition holds for the maps
$\text{id} : U_{j_1j_2} \to U_{j_2j_1}$.
Ok, so now we can use Schemes, Lemma \href{https://stacks.math.columbia.edu/tag/01JC}{lemma glue schemes (Tag 01JC)}
to glue the system
$((U_{j, i_5})_j, (W_{j_1, j_2, i_5})_{j_1, j_2},
(\varphi_{j_1, j_2, i_5})_{j_1, j_2})$ to get a scheme
$X_{i_5} \to S_{i_5}$. By construction the base change of
$X_{i_5}$ to $S$ is formed by glueing the open affines
$U_j$ along the opens $U_{j_1} \leftarrow U_{j_1j_2} \rightarrow U_{j_2}$.
Hence $S \times_{S_{i_5}} X_{i_5} \cong X$ as desired.
\end{proof}


\begin{lemma}
\label{lemma-locally-finite-type-in-finite-presentation}
Let $f : X \to S$ be a morphism of schemes.
Assume:
\begin{enumerate}
\item The morphism $f$ is locally of finite type.
\item The scheme $X$ is quasi-compact and quasi-separated.
\end{enumerate}
Then there exists a morphism of finite presentation
$f' : X' \to S$ and an immersion $X \to X'$ of schemes over $S$.
\end{lemma}

\begin{proof}
By
Proposition \ref{proposition-approximate}
we can write
$X = \lim_i X_i$ with each $X_i$ of finite type over $\mathbf{Z}$ and
with transition morphisms $f_{ii'} : X_i \to X_{i'}$ affine.
Consider the commutative diagram
$$
\xymatrix{
X \ar[r] \ar[rd] & X_{i, S} \ar[r] \ar[d] & X_i \ar[d] \\
& S \ar[r] & \Spec(\mathbf{Z})
}
$$
Note that $X_i$ is of finite presentation over $\Spec(\mathbf{Z})$, see
Morphisms,
Lemma \href{https://stacks.math.columbia.edu/tag/01TX}{lemma noetherian finite type finite presentation (Tag 01TX)}.
Hence the base change $X_{i, S} \to S$ is of finite presentation by
Morphisms, Lemma \href{https://stacks.math.columbia.edu/tag/01TS}{lemma base change finite presentation (Tag 01TS)}.
Thus it suffices to show that the arrow $X \to X_{i, S}$ is an
immersion for $i$ sufficiently large.

\medskip\noindent
To do this we choose a finite affine open covering
$X = V_1 \cup \ldots \cup V_n$ such that
$f$ maps each $V_j$ into an affine open $U_j \subset S$.
Let $h_{j, a} \in \mathcal{O}_X(V_j)$ be a finite
set of elements which generate $\mathcal{O}_X(V_j)$ as
an $\mathcal{O}_S(U_j)$-algebra, see
Morphisms, Lemma \href{https://stacks.math.columbia.edu/tag/01T2}{lemma locally finite type characterize (Tag 01T2)}.
By Lemmas \ref{lemma-descend-opens} and \ref{lemma-limit-affine}
(after possibly shrinking $I$) we may assume that
there exist affine open coverings
$X_i = V_{1, i} \cup \ldots \cup V_{n, i}$
compatible with transition maps such that $V_j = \lim_i V_{j, i}$.
By Lemma \ref{lemma-descend-section} we can choose $i$ so large that each
$h_{j, a}$ comes from an element
$h_{j, a, i} \in \mathcal{O}_{X_i}(V_{j, i})$.
Thus the arrow in
$$
V_j \longrightarrow U_j \times_{\Spec(\mathbf{Z})} V_{j, i} =
(V_{j, i})_{U_j} \subset (V_{j, i})_S \subset X_{i, S}
$$
is a closed immersion. Since $\bigcup (V_{j, i})_{U_j}$ forms an open of
$X_{i, S}$ and since the inverse image of $(V_{j, i})_{U_j}$ in $X$
is $V_j$ it follows that $X \to X_{i, S}$ is an immersion.
\end{proof}


\begin{lemma}
\label{lemma-finite-type-closed-in-finite-presentation}
Let $f : X \to S$ be a morphism of schemes.
Assume:
\begin{enumerate}
\item The morphism $f$ is of locally of finite type.
\item The scheme $X$ is quasi-compact and quasi-separated, and
\item The scheme $S$ is quasi-separated.
\end{enumerate}
Then there exists a morphism of finite presentation
$f' : X' \to S$ and a closed immersion $X \to X'$ of schemes over $S$.
\end{lemma}

\begin{proof}
By Lemma \ref{lemma-locally-finite-type-in-finite-presentation} above
there exists a morphism $Y \to S$ of finite presentation and an
immersion $i : X \to Y$ of schemes over $S$.
For every point $x \in X$, there exists an affine open
$V_x \subset Y$ such that $i^{-1}(V_x) \to V_x$ is a
closed immersion. Since $X$ is quasi-compact we can find
finitely may affine opens $V_1, \ldots, V_n \subset Y$
such that $i(X) \subset V_1 \cup \ldots \cup V_n$ and
$i^{-1}(V_j) \to V_j$ is a closed immersion. In other words
such that $i : X \to X' = V_1 \cup \ldots \cup V_n$ is a
closed immersion of schemes over $S$.
Since $S$ is quasi-separated and $Y$ is quasi-separated over $S$
we deduce that $Y$ is quasi-separated, see
Schemes, Lemma \href{https://stacks.math.columbia.edu/tag/01KU}{lemma separated permanence (Tag 01KU)}.
Hence the open immersion $X' = V_1 \cup \ldots \cup V_n \to Y$
is quasi-compact. This implies that
$X' \to Y$ is of finite presentation, see
Morphisms,
Lemma \href{https://stacks.math.columbia.edu/tag/01TU}{lemma quasi compact open immersion finite presentation (Tag 01TU)}.
We conclude since then $X' \to Y \to S$ is a composition of morphisms
of finite presentation, and hence of finite presentation (see
Morphisms, Lemma \href{https://stacks.math.columbia.edu/tag/01TR}{lemma composition finite presentation (Tag 01TR)}).
\end{proof}


\begin{lemma}
\label{lemma-finite-type-is-limit-finite-presentation}
Let $f : X \to S$ be a morphism of schemes. Assume
\begin{enumerate}
\item The morphism $f$ is of locally of finite type.
\item The scheme $X$ is quasi-compact and quasi-separated, and
\item The scheme $S$ is quasi-separated.
\end{enumerate}
Then $X = \lim X_i$ where the $X_i \to S$ are of
finite presentation, the $X_i$ are quasi-compact and quasi-separated,
and the transition morphisms $X_{i'} \to X_i$ are closed immersions
(which implies that $X \to X_i$ are closed immersions for all $i$).
\end{lemma}

\begin{proof}
By Lemma \ref{lemma-finite-type-closed-in-finite-presentation}
there is a closed immersion $X \to Y$ with $Y \to S$ of
finite presentation. Then $Y$ is quasi-separated by
Schemes, Lemma \href{https://stacks.math.columbia.edu/tag/01KU}{lemma separated permanence (Tag 01KU)}.
Since $X$ is quasi-compact, we may assume
$Y$ is quasi-compact by replacing $Y$ with a quasi-compact open
containing $X$. We see that $X = \lim X_i$ with $X_i \to Y$ a closed
immersion of finite presentation by
Lemma \ref{lemma-closed-is-limit-closed-and-finite-presentation}.
The morphisms $X_i \to S$ are of finite presentation by
Morphisms, Lemma \href{https://stacks.math.columbia.edu/tag/01TR}{lemma composition finite presentation (Tag 01TR)}.
\end{proof}


\begin{lemma}
\label{lemma-closed-is-limit-closed-and-finite-presentation}
Let $X \to Y$ be a closed immersion of schemes. Assume $Y$ quasi-compact and
quasi-separated. Then $X$ can be written as a directed limit $X = \lim X_i$
of schemes over $Y$ where $X_i \to Y$ is a closed immersion
of finite presentation.
\end{lemma}

\begin{proof}
Let $\mathcal{I} \subset \mathcal{O}_Y$ be the quasi-coherent sheaf of
ideals defining $X$ as a closed subscheme of $Y$. By
Properties, Lemma \href{https://stacks.math.columbia.edu/tag/01PG}{lemma quasi coherent colimit finite type (Tag 01PG)}
we can write $\mathcal{I}$ as a directed colimit
$\mathcal{I} = \colim_{i \in I} \mathcal{I}_i$ of its
quasi-coherent sheaves of ideals of finite type.
Let $X_i \subset Y$ be the closed subscheme defined by $\mathcal{I}_i$.
These form an inverse system of schemes indexed by $I$.
The transition morphisms $X_i \to X_{i'}$ are affine because
they are closed immersions. Each $X_i$ is quasi-compact and quasi-separated
since it is a closed subscheme of $Y$ and $Y$ is quasi-compact and
quasi-separated by our assumptions.
We have $X = \lim_i X_i$ as follows directly from the
fact that $\mathcal{I} = \colim_{i \in I} \mathcal{I}_a$.
Each of the morphisms $X_i \to Y$ is of finite presentation, see
Morphisms, Lemma \href{https://stacks.math.columbia.edu/tag/01TV}{lemma closed immersion finite presentation (Tag 01TV)}.
\end{proof}


\begin{lemma}
\label{lemma-eventually-separated}
Let $S$ be a scheme. Let $X = \lim X_i$ be a directed
limit of schemes over $S$ with affine transition morphisms.
Assume
\begin{enumerate}
\item $S$ quasi-separated,
\item $X_i$ quasi-compact and quasi-separated,
\item $X \to S$ separated.
\end{enumerate}
Then $X_i \to S$ is separated for all $i$ large enough.
\end{lemma}

\begin{proof}
Let $0 \in I$. Note that $I$ is nonempty as the limit is directed.
As $X_0$ is quasi-compact we can find finitely many
affine opens $U_1, \ldots, U_n \subset S$ such that
$X_0 \to S$ maps into $U_1 \cup \ldots \cup U_n$.
Denote $h_i : X_i \to S$ the structure morphism.
It suffices to check that for some $i \geq 0$ the morphisms
$h_i^{-1}(U_j) \to U_j$ are separated for $j = 1, \ldots,  n$.
Since $S$ is quasi-separated the morphisms $U_j \to S$ are quasi-compact.
Hence $h_i^{-1}(U_j)$ is quasi-compact and quasi-separated.
In this way we reduce to the case $S$ affine. In this case we
have to show that $X_i$ is separated and we know that $X$ is separated.
Thus the lemma follows from Lemma \ref{lemma-limit-separated}.
\end{proof}


\begin{lemma}
\label{lemma-limit-separated}
In Situation \ref{situation-descent} if $S$ is separated,
then for some $i_0 \in I$ the schemes $S_i$ for $i \geq i_0$
are separated.
\end{lemma}

\begin{proof}
Choose a finite affine open covering
$S_0 = U_{0, 1} \cup \ldots \cup U_{0, m}$.
Set $U_{i, j} \subset S_i$ and $U_j \subset S$
equal to the inverse image of $U_{0, j}$.
Note that $U_{i, j}$ and $U_j$ are affine. As $S$ is separated
the intersections $U_{j_1} \cap U_{j_2}$ are affine. Since
$U_{j_1} \cap U_{j_2} = \lim_{i \geq 0} U_{i, j_1} \cap U_{i, j_2}$
we see that $U_{i, j_1} \cap U_{i, j_2}$ is affine for large $i$
by Lemma \ref{lemma-limit-affine}. To show that $S_i$ is separated
for large $i$ it now suffices to show that
$$
\mathcal{O}_{S_i}(U_{i, j_1})
\otimes_{\mathcal{O}_S(S)}
\mathcal{O}_{S_i}(U_{i, j_2})
\longrightarrow
\mathcal{O}_{S_i}(U_{i, j_1} \cap U_{i, j_2})
$$
is surjective for large $i$
(Schemes, Lemma \href{https://stacks.math.columbia.edu/tag/01KP}{lemma characterize separated (Tag 01KP)}).

\medskip\noindent
To get rid of the annoying indices, assume we have affine opens
$U, V \subset S_0$ such that $U \cap V$ is affine too.
Let $U_i, V_i \subset S_i$, resp.\ $U, V \subset S$ be the inverse images.
We have to show that
$\mathcal{O}(U_i) \otimes \mathcal{O}(V_i) \to
\mathcal{O}(U_i \cap V_i)$
is surjective for $i$ large enough and we know that
$\mathcal{O}(U) \otimes \mathcal{O}(V) \to \mathcal{O}(U \cap V)$
is surjective. Note that
$\mathcal{O}(U_0) \otimes \mathcal{O}(V_0) \to
\mathcal{O}(U_0 \cap V_0)$
is of finite type, as the diagonal morphism $S_i \to S_i \times S_i$
is an immersion (Schemes, Lemma \href{https://stacks.math.columbia.edu/tag/01KJ}{lemma diagonal immersion (Tag 01KJ)})
hence locally of finite type
(Morphisms, Lemmas \href{https://stacks.math.columbia.edu/tag/01T2}{lemma locally finite type characterize (Tag 01T2)} and
\href{https://stacks.math.columbia.edu/tag/01T5}{lemma immersion locally finite type (Tag 01T5)}).
Thus we can choose elements
$f_{0, 1}, \ldots, f_{0, n} \in \mathcal{O}(U_0 \cap V_0)$
which generate $\mathcal{O}(U_0 \cap V_0)$ over
$\mathcal{O}(U_0) \otimes \mathcal{O}(V_0)$.
Observe that for $i \geq 0$ the diagram of schemes
$$
\xymatrix{
U_i \cap V_i \ar[r] \ar[d] & U_i \ar[d] \\
U_0 \cap V_0 \ar[r] & U_0
}
$$
is cartesian. Thus we see that the images
$f_{i, 1}, \ldots, f_{i, n} \in \mathcal{O}(U_i \cap V_i)$
generate $\mathcal{O}(U_i \cap V_i)$ over
$\mathcal{O}(U_i) \otimes \mathcal{O}(V_0)$
and a fortiori over
$\mathcal{O}(U_i) \otimes \mathcal{O}(V_i)$.
By assumption the images $f_1, \ldots, f_n \in \mathcal{O}(U \otimes V)$
are in the image of the map
$\mathcal{O}(U) \otimes \mathcal{O}(V) \to \mathcal{O}(U \cap V)$.
Since
$\mathcal{O}(U) \otimes \mathcal{O}(V) =
\colim \mathcal{O}(U_i) \otimes \mathcal{O}(V_i)$
we see that they are in the image of the map at some finite level
and the lemma is proved.
\end{proof}


\begin{lemma}
\label{lemma-descend-separated-finite-presentation}
Notation and assumptions as in Situation \ref{situation-descent-property}.
If $f$ is separated, then $f_i$ is separated for some $i \geq 0$.
\end{lemma}

\begin{proof}
Apply Lemma \ref{lemma-descend-closed-immersion-finite-presentation}
to the diagonal morphism $\Delta_{X_0/S_0} : X_0 \to X_0 \times_{S_0} X_0$.
(This is permissible as diagonal morphisms are locally of finite type
and the fibre product $X_0 \times_{S_0} X_0$ is quasi-compact and
quasi-separated, see
Schemes, Lemma \href{https://stacks.math.columbia.edu/tag/01KJ}{lemma diagonal immersion (Tag 01KJ)},
Morphisms, Lemma \href{https://stacks.math.columbia.edu/tag/01T5}{lemma immersion locally finite type (Tag 01T5)}, and
Schemes, Remark \href{https://stacks.math.columbia.edu/tag/0816}{remark quasi compact and quasi separated (Tag 0816)}.
\end{proof}


\begin{lemma}
\label{lemma-descend-closed-immersion-finite-presentation}
Notation and assumptions as in Situation \ref{situation-descent-property}.
If
\begin{enumerate}
\item $f$ is a closed immersion, and
\item $f_0$ is locally of finite type,
\end{enumerate}
then there exists an $i \geq 0$ such that $f_i$ is a closed immersion.
\end{lemma}

\begin{proof}
A closed immersion is affine, see
Morphisms, Lemma \href{https://stacks.math.columbia.edu/tag/01SE}{lemma closed immersion affine (Tag 01SE)}.
Hence by Lemma \ref{lemma-descend-affine-finite-presentation} above
after increasing $0$ we may assume that $f_0$ is affine.
By writing $Y_0$ as a finite union of affines we reduce to proving
the result when $X_0$ and $Y_0$ are affine and map
into a common affine $W \subset S_0$. The corresponding algebra
statement is a consequence of
Algebra, Lemma \href{https://stacks.math.columbia.edu/tag/07RH}{lemma colimit surjective (Tag 07RH)}.
\end{proof}


\begin{lemma}
\label{lemma-descend-affine-finite-presentation}
Notation and assumptions as in Situation \ref{situation-descent-property}.
If $f$ is affine, then there exists an index $i \geq 0$
such that $f_i$ is affine.
\end{lemma}

\begin{proof}
Let $Y_0 = \bigcup_{j = 1, \ldots, m} V_{j, 0}$ be a finite affine
open covering. Set $U_{j, 0} = f_0^{-1}(V_{j, 0})$. For $i \geq 0$
we denote $V_{j, i}$ the inverse image of $V_{j, 0}$ in $Y_i$ and
$U_{j, i} = f_i^{-1}(V_{j, i})$. Similarly we have
$U_j = f^{-1}(V_j)$. Then $U_j = \lim_{i \geq 0} U_{j, i}$
(see Lemma \ref{lemma-directed-inverse-system-has-limit}).
Since $U_j$ is affine by assumption we see that
each $U_{j, i}$ is affine for $i$ large enough, see
Lemma \ref{lemma-limit-affine}. As there are finitely many $j$ we
can pick an $i$ which works for all $j$. Thus $f_i$ is
affine for $i$ large enough, see
Morphisms, Lemma \href{https://stacks.math.columbia.edu/tag/01S8}{lemma characterize affine (Tag 01S8)}.
\end{proof}


\begin{lemma}
\label{lemma-descend-opens}
In Situation \ref{situation-descent} we have the following:
\begin{enumerate}
\item Given any quasi-compact open $V \subset S = \lim_i S_i$
there exists an $i \in I$ and a quasi-compact open $V_i \subset S_i$
such that $f_i^{-1}(V_i) = V$.
\item Given $V_i \subset S_i$ and $V_{i'} \subset S_{i'}$
quasi-compact opens such that $f_i^{-1}(V_i) = f_{i'}^{-1}(V_{i'})$
there exists an index $i'' \geq i, i'$ such that
$f_{i''i}^{-1}(V_i) = f_{i''i'}^{-1}(V_{i'})$.
\item If $V_{1, i}, \ldots, V_{n, i} \subset S_i$ are quasi-compact
opens and $S = f_i^{-1}(V_{1, i}) \cup \ldots \cup f_i^{-1}(V_{n, i})$
then $S_{i'} = f_{i'i}^{-1}(V_{1, i}) \cup \ldots \cup f_{i'i}^{-1}(V_{n, i})$
for some $i' \geq i$.
\end{enumerate}
\end{lemma}

\begin{proof}
Choose $i_0 \in I$. Note that $I$ is nonempty as the limit is directed.
For convenience we write $S_0 = S_{i_0}$ and $i_0 = 0$.
Choose an affine open covering $S_0 = U_{1, 0} \cup \ldots \cup U_{m, 0}$.
Denote $U_{j, i} \subset S_i$ the inverse image of $U_{j, 0}$
under the transition morphism for $i \geq 0$.
Denote $U_j$ the inverse image of $U_{j, 0}$ in $S$.
Note that $U_j = \lim_i U_{j, i}$ is a limit of affine
schemes.

\medskip\noindent
We first prove the uniqueness statement: Let
$V_i \subset S_i$ and $V_{i'} \subset S_{i'}$
quasi-compact opens such that $f_i^{-1}(V_i) = f_{i'}^{-1}(V_{i'})$.
It suffices to show that $f_{i''i}^{-1}(V_i \cap U_{j, i''})$ and
$f_{i''i'}^{-1}(V_{i'} \cap U_{j, i''})$ become equal
for $i''$ large enough. Hence we reduce to the case
of a limit of affine schemes. In this case write
$S = \Spec(R)$ and $S_i = \Spec(R_i)$ for all $i \in I$.
We may write $V_i = S_i \setminus V(h_1, \ldots, h_m)$
and $V_{i'} = S_{i'} \setminus V(g_1, \ldots, g_n)$.
The assumption means that the ideals
$\sum g_jR$ and $\sum h_jR$ have the same radical
in $R$. This means that $g_j^N = \sum a_{jj'}h_{j'}$ and
$h_j^N = \sum b_{jj'} g_{j'}$ for some $N \gg 0$ and $a_{jj'}$
and $b_{jj'}$ in $R$.
Since $R = \colim_i R_i$ we can chose an index
$i'' \geq i$ such that the equations
$g_j^N = \sum a_{jj'}h_{j'}$ and
$h_j^N = \sum b_{jj'} g_{j'}$ hold in $R_{i''}$ for some
$a_{jj'}$ and $b_{jj'}$ in $R_{i''}$. This implies that
the ideals $\sum g_jR_{i''}$ and $\sum h_jR_{i''}$ have the same radical
in $R_{i''}$ as desired.

\medskip\noindent
We prove existence: If $S_0$ is affine, then $S_i = \Spec(R_i)$ for all
$i \geq 0$ and $S = \Spec(R)$ with $R = \colim R_i$. Then
$V = S \setminus V(g_1, \ldots, g_n)$ for some $g_1, \ldots, g_n \in R$.
Choose any $i$ large enough so that each of the $g_j$ comes from an
element $g_{j, i} \in R_i$ and take
$V_i = S_i \setminus V(g_{1, i}, \ldots, g_{n, i})$.
If $S_0$ is general, then the opens $V \cap U_j$
are quasi-compact because $S$ is quasi-separated. Hence by the
affine case we see that for each $j = 1, \ldots, m$
there exists an $i_j \in I$ and a quasi-compact open
$V_{i_j} \subset U_{j, i_j}$ whose inverse image in $U_j$
is $V \cap U_j$. Set $i = \max(i_1, \ldots, i_m)$
and let $V_i = \bigcup f_{ii_j}^{-1}(V_{i_j})$.

\medskip\noindent
The statement on coverings follows from the uniqueness statement
for the opens $V_{1, i} \cup \ldots \cup V_{n, i}$ and $S_i$ of $S_i$.
\end{proof}


\begin{lemma}
\label{lemma-limit-quasi-affine}
In Situation \ref{situation-descent} if $S$ is quasi-affine, then
for some $i_0 \in I$ the schemes $S_i$ for $i \geq i_0$ are quasi-affine.
\end{lemma}

\begin{proof}
Choose $i_0 \in I$. Note that $I$ is nonempty as the limit is directed.
For convenience we write $S_0 = S_{i_0}$ and $i_0 = 0$.
Let $s \in S$. We may choose an affine open
$U_0 \subset S_0$ containing $f_0(s)$. Since $S$ is quasi-affine
we may choose an element $a \in \Gamma(S, \mathcal{O}_S)$ such
that $s \in D(a) \subset f_0^{-1}(U_0)$, and such that
$D(a)$ is affine. By Lemma \ref{lemma-descend-section}
there exists an $i \geq 0$ such that $a$
comes from an element $a_i \in \Gamma(S_i, \mathcal{O}_{S_i})$.
For any index $j \geq i$ we denote $a_j$
the image of $a_i$ in the global sections of the
structure sheaf of $S_j$.
Consider the opens $D(a_j) \subset S_j$
and $U_j = f_{j0}^{-1}(U_0)$. Note that
$U_j$ is affine and $D(a_j)$ is a quasi-compact open of $S_j$,
see Properties, Lemma \href{https://stacks.math.columbia.edu/tag/01PV}{lemma affine cap s open (Tag 01PV)}
for example. Hence we may apply Lemma \ref{lemma-descend-opens} to the opens
$U_j$ and $U_j \cup D(a_j)$ to conclude that
$D(a_j) \subset U_j$ for some  $j \geq i$.
For such an index $j$ we see that $D(a_j) \subset S_j$ is an affine open
(because $D(a_j)$ is a standard affine open of the affine open $U_j$)
containing the image $f_j(s)$.

\medskip\noindent
We conclude that for every $s \in S$ there exist
an index $i \in I$, and a global section
$a \in \Gamma(S_i, \mathcal{O}_{S_i})$
such that $D(a) \subset S_i$ is an affine open
containing $f_i(s)$. Because $S$ is quasi-compact we
may choose a single index $i \in I$ and global sections
$a_1, \ldots, a_m \in \Gamma(S_i, \mathcal{O}_{S_i})$
such that each $D(a_j) \subset S_i$ is affine open
and such that $f_i : S \to S_i$ has image contained
in the union $W_i = \bigcup_{j = 1, \ldots, m} D(a_j)$.
For $i' \geq i$ set $W_{i'} = f_{i'i}^{-1}(W_i)$.
Since $f_i^{-1}(W_i)$ is all of $S$ we see
(by Lemma \ref{lemma-descend-opens} again)
that for a suitable $i' \geq i$ we
have $S_{i'} = W_{i'}$. Thus we may replace $i$ by
$i'$ and assume that $S_i = \bigcup_{j = 1, \ldots, m} D(a_j)$.
This implies that $\mathcal{O}_{S_i}$ is an ample invertible
sheaf on $S_i$ (see Properties, Definition \ref{properties-definition-ample})
and hence that $S_i$ is quasi-affine, see
Properties, Lemma \href{https://stacks.math.columbia.edu/tag/01QE}{lemma quasi affine O ample (Tag 01QE)}.
Hence we win.
\end{proof}


\begin{lemma}
\label{lemma-descend-section}
In Situation \ref{situation-descent}.
Suppose that $\mathcal{F}_0$ is a quasi-coherent sheaf on $S_0$.
Set $\mathcal{F}_i = f_{i0}^*\mathcal{F}_0$ for $i \geq 0$ and set
$\mathcal{F} = f_0^*\mathcal{F}_0$.
Then
$$
\Gamma(S, \mathcal{F}) = \colim_{i \geq 0} \Gamma(S_i, \mathcal{F}_i)
$$
\end{lemma}

\begin{proof}
Write $\mathcal{A}_j = f_{i0, *} \mathcal{O}_{S_i}$.
This is a quasi-coherent sheaf of $\mathcal{O}_{S_0}$-algebras
(see Morphisms, Lemma \href{https://stacks.math.columbia.edu/tag/01SA}{lemma affine equivalence algebras (Tag 01SA)})
and $S_i$ is the relative spectrum of $\mathcal{A}_i$ over $S_0$.
In the proof of Lemma \ref{lemma-directed-inverse-system-has-limit}
we constructed $S$ as the relative spectrum of
$\mathcal{A} = \colim_{i \geq 0} \mathcal{A}_i$
over $S_0$. Set
$$
\mathcal{M}_i = \mathcal{F}_0 \otimes_{\mathcal{O}_{S_0}} \mathcal{A}_i
$$
and
$$
\mathcal{M} = \mathcal{F}_0 \otimes_{\mathcal{O}_{S_0}} \mathcal{A}.
$$
Then we have $f_{i0, *} \mathcal{F}_i = \mathcal{M}_i$
and $f_{0, *}\mathcal{F} = \mathcal{M}$. Since $\mathcal{A}$
is the colimit of the sheaves $\mathcal{A}_i$ and since tensor
product commutes with directed colimits, we conclude that
$\mathcal{M} = \colim_{i \geq 0} \mathcal{M}_i$.
Since $S_0$ is quasi-compact and quasi-separated we see that
\begin{eqnarray*}
\Gamma(S, \mathcal{F})
& = &
\Gamma(S_0, \mathcal{M}) \\
& = &
\Gamma(S_0, \colim_{i \geq 0} \mathcal{M}_i) \\
& = &
\colim_{i \geq 0} \Gamma(S_0, \mathcal{M}_i) \\
& = &
\colim_{i \geq 0} \Gamma(S_i, \mathcal{F}_i)
\end{eqnarray*}
see Sheaves, Lemma \href{https://stacks.math.columbia.edu/tag/009F}{lemma directed colimits sections (Tag 009F)} and
Topology, Lemma \href{https://stacks.math.columbia.edu/tag/0069}{lemma topology quasi separated scheme (Tag 0069)}
for the middle equality.
\end{proof}


\begin{lemma}
\label{lemma-quasi-affine-finite-type-over-Z}
Let $W$ be a quasi-affine scheme of finite type over
$\mathbf{Z}$. Suppose $W \to \Spec(R)$ is an
open immersion into an affine scheme. There exists a
finite type $\mathbf{Z}$-algebra $A \subset R$
which induces an open immersion $W \to \Spec(A)$.
Moreover, $R$ is the directed colimit of such subalgebras.
\end{lemma}

\begin{proof}
Choose an affine open covering $W = \bigcup_{i = 1, \ldots, n} W_i$
such that each $W_i$ is a standard affine open in $\Spec(R)$.
In other words, if we write $W_i = \Spec(R_i)$
then $R_i = R_{f_i}$ for some $f_i \in R$.
Choose finitely many $x_{ij} \in R_i$ which generate
$R_i$ over $\mathbf{Z}$.
Pick an $N \gg 0$ such that each $f_i^Nx_{ij}$ comes from an
element of $R$, say $y_{ij} \in R$.
Set $A$ equal to the $\mathbf{Z}$-algebra generated by
the $f_i$ and the $y_{ij}$ and (optionally) finitely many
additional elements of $R$. Then $A$ works. Details omitted.
\end{proof}


\begin{lemma}
\label{lemma-diagram}
Suppose given a cartesian diagram of rings
$$
\xymatrix{
B \ar[r]_s & R \\
B'\ar[u] \ar[r] & R' \ar[u]_t
}
$$
Let $W' \subset \Spec(R')$ be an open of
the form $W' = D(f_1) \cup \ldots \cup D(f_n)$
such that $t(f_i) = s(g_i)$ for some $g_i \in B$
and $B_{g_i} \cong R_{s(g_i)}$. Then $B' \to R'$
induces an open immersion of $W'$ into $\Spec(B')$.
\end{lemma}

\begin{proof}
Set $h_i = (g_i, f_i) \in B'$. More on Algebra,
Lemma \ref{more-algebra-lemma-diagram-localize} shows that
$(B')_{h_i} \cong (R')_{f_i}$ as desired.
\end{proof}


\begin{lemma}
\label{more-algebra-lemma-diagram-localize}
Suppose given a cartesian diagram of rings
$$
\xymatrix{
R &
R' \ar[l]^t \\
B \ar[u]_s &
B'\ar[u] \ar[l]
}
$$
i.e., $B' = B \times_R R'$. If $h \in B'$ corresponds to $g \in B$
and $f \in R'$ such that $s(g) = t(f)$, then the diagram
$$
\xymatrix{
R_{s(g)} = R_{t(f)} &
(R')_f \ar[l]^-t \\
B_g \ar[u]_s &
(B')_h \ar[u] \ar[l]
}
$$
is cartesian too.
\end{lemma}

\begin{proof}
The equality $B' = B \times_R R'$ tells us that
$$
0 \to B' \to B \oplus R' \xrightarrow{s, -t} R
$$
is an exact sequence of $B'$-modules. We have $B_g = B_h$,
$R'_f = R'_h$, and $R_{s(g)} = R_{t(f)} = R_h$ as $B'$-modules.
By exactness of localization
(Algebra, Proposition \href{https://stacks.math.columbia.edu/tag/00CS}{proposition localization exact (Tag 00CS)})
we find that
$$
0 \to B'_h \to B_g \oplus R'_f \xrightarrow{s, -t} R_{s(g)} = R_{t(f)}
$$
is an exact sequence. This proves the lemma.
\end{proof}


\begin{lemma}
\label{lemma-directed-inverse-system-has-limit}
Let $I$ be a directed set. Let $(S_i, f_{ii'})$ be an
inverse system of schemes over $I$. If all the morphisms
$f_{ii'} : S_i \to S_{i'}$ are affine, then the limit $S = \lim_i S_i$ exists
in the category of schemes. Moreover,
\begin{enumerate}
\item each of the morphisms $f_i : S \to S_i$ is affine,
\item for an element $0 \in I$ and any open subscheme $U_0 \subset S_0$
we have
$$
f_0^{-1}(U_0) = \lim_{i \geq 0} f_{i0}^{-1}(U_0)
$$
in the category of schemes.
\end{enumerate}
\end{lemma}

\begin{proof}
Choose an element $0 \in I$. Note that $I$ is nonempty as the limit is
directed. For every $i \geq 0$ consider the quasi-coherent sheaf of
$\mathcal{O}_{S_0}$-algebras $\mathcal{A}_i = f_{i0, *}\mathcal{O}_{S_i}$.
Recall that $S_i = \underline{\Spec}_{S_0}(\mathcal{A}_i)$,
see Morphisms, Lemma \href{https://stacks.math.columbia.edu/tag/01S8}{lemma characterize affine (Tag 01S8)}.
Set $\mathcal{A} = \colim_{i \geq 0} \mathcal{A}_i$.
This is a quasi-coherent sheaf of $\mathcal{O}_{S_0}$-algebras,
see Schemes, Section \href{https://stacks.math.columbia.edu/tag/01LA}{section quasi coherent (Tag 01LA)}.
Set $S = \underline{\Spec}_{S_0}(\mathcal{A})$.
By Morphisms, Lemma \href{https://stacks.math.columbia.edu/tag/01SA}{lemma affine equivalence algebras (Tag 01SA)}
we get for $i \geq 0$ morphisms $f_i : S \to S_i$ compatible with
the transition morphisms. Note that the morphisms $f_i$ are
affine by Morphisms, Lemma \href{https://stacks.math.columbia.edu/tag/01SG}{lemma affine permanence (Tag 01SG)} for example.
By Lemma \href{https://stacks.math.columbia.edu/tag/01YW}{lemma directed inverse system affine schemes has limit (Tag 01YW)} above
we see that for any affine open $U_0 \subset S_0$ the
inverse image $U = f_0^{-1}(U_0) \subset S$ is the limit of the
system of opens $U_i = f_{i0}^{-1}(U_0)$, $i \geq 0$ in the
category of schemes.

\medskip\noindent
Let $T$ be a scheme. Let $g_i : T \to S_i$ be a compatible system
of morphisms. To show that $S = \lim_i S_i$ we have
to prove there is a unique morphism $g : T \to S$ with
$g_i = f_i \circ g$ for all $i \in I$.
For every $t \in T$ there exists an affine open
$U_0 \subset S_0$ containing $g_0(t)$. Let $V \subset g_0^{-1}(U_0)$
be an affine open neighbourhood containing $t$.
By the remarks above we obtain a unique morphism
$g_V : V \to U = f_0^{-1}(U_0)$ such that $f_i \circ g_V = g_i|_{U_i}$
for all $i$. The open sets $V \subset T$ so constructed form
a basis for the topology of $T$. The morphisms $g_V$ glue to a morphism
$g : T \to S$ because of the uniqueness property. This gives the
desired morphism $g : T \to S$.

\medskip\noindent
The final statement is clear from the construction of the limit above.
\end{proof}


\begin{lemma}
\label{lemma-limit-affine}
In Situation \ref{situation-descent} if $S$ is affine,
then for some $i_0 \in I$ the schemes $S_i$ for $i \geq i_0$
are affine.
\end{lemma}

\begin{proof}
By Lemma \ref{lemma-limit-quasi-affine} we may assume that $S_0$ is
quasi-affine for some $0 \in I$. Set $R_0 = \Gamma(S_0, \mathcal{O}_{S_0})$.
Then $S_0$ is a quasi-compact open of $T_0 = \Spec(R_0)$. Denote
$j_0 : S_0 \to T_0$ the corresponding quasi-compact open immersion.
For $i \geq 0$ set $\mathcal{A}_i = f_{i0, *}\mathcal{O}_{S_i}$.
Since $f_{i0}$ is affine we see that
$S_i = \underline{\Spec}_{S_0}(\mathcal{A}_i)$.
Set $T_i = \underline{\Spec}_{T_0}(j_{0, *}\mathcal{A}_i)$.
Then $T_i \to T_0$ is affine, hence $T_i$ is affine. Thus
$T_i$ is the spectrum of
$$
R_i = \Gamma(T_0, j_{0, *}\mathcal{A}_i) = \Gamma(S_0, \mathcal{A}_i) =
\Gamma(S_i, \mathcal{O}_{S_i}).
$$
Write $S = \Spec(R)$. We have $R = \colim_i R_i$
by Lemma \ref{lemma-descend-section}.
Hence also $S = \lim_i T_i$. As formation of the relative spectrum commutes
with base change, the inverse image
of the open $S_0 \subset T_0$ in $T_i$ is $S_i$.
Let $Z_0 = T_0 \setminus S_0$ and let $Z_i \subset T_i$
be the inverse image of $Z_0$. As $S_i = T_i \setminus Z_i$, it suffices
to show that $Z_i$ is empty for some $i$. Assume $Z_i$ is nonempty for all
$i$ to get a contradiction. By Lemma \ref{lemma-limit-closed-nonempty}
there exists a point $s$ of $S = \lim T_i$ which maps to a point of $Z_i$
for every $i$. But $S = \lim_i S_i$, and hence we arrive at a contradiction
by Lemma \ref{lemma-topology-limit}.
\end{proof}


\begin{lemma}
\label{lemma-limit-closed-nonempty}
In Situation \ref{situation-descent}.
Suppose for each $i$ we are given a nonempty closed subset
$Z_i \subset S_i$ with $f_{i'i}(Z_{i'}) \subset Z_i$ for all
$i' \geq i$.
Then there exists a point $s \in S$ with $f_i(s) \in Z_i$ for
all $i$.
\end{lemma}

\begin{proof}
Let $Z_i \subset S_i$ also denote the reduced closed subscheme
associated to $Z_i$, see Schemes,
Definition \ref{schemes-definition-reduced-induced-scheme}.
A closed immersion is affine, and a composition of affine
morphisms is affine (see
Morphisms, Lemmas \href{https://stacks.math.columbia.edu/tag/01SE}{lemma closed immersion affine (Tag 01SE)}
and \href{https://stacks.math.columbia.edu/tag/01SC}{lemma composition affine (Tag 01SC)}), and hence $Z_{i'} \to S_i$ is
affine when $i' \geq i$. We conclude that the morphism
$f_{i'i} : Z_{i'} \to Z_i$ is affine by
Morphisms, Lemma \href{https://stacks.math.columbia.edu/tag/01SG}{lemma affine permanence (Tag 01SG)}.
Each of the schemes $Z_i$ is quasi-compact as a closed
subscheme of a quasi-compact scheme. Hence we may apply
Lemma \ref{lemma-limit-nonempty} to see that
$Z = \lim_i Z_i$ is nonempty. Since there is a
canonical morphism $Z \to S$ we win.
\end{proof}


\begin{lemma}
\label{lemma-limit-nonempty}
Let $S = \lim S_i$ be the limit of a directed inverse system
of schemes with affine transition morphisms
(Lemma \ref{lemma-directed-inverse-system-has-limit}).
If all the schemes $S_i$ are nonempty and quasi-compact,
then the limit $S = \lim_i S_i$ is nonempty.
\end{lemma}

\begin{proof}
Choose $0 \in I$. Note that $I$ is nonempty as the limit is directed.
Choose an affine open covering $S_0 = \bigcup_{j = 1, \ldots, m} U_j$.
Since $I$ is directed there exists a $j \in \{1, \ldots, m\}$
such that $f_{i0}^{-1}(U_j) \not = \emptyset$ for all
$i \geq 0$. Hence $\lim_{i \geq 0} f_{i0}^{-1}(U_j)$ is not
empty since a directed colimit of nonzero rings is nonzero
(because $1 \not = 0$). As $\lim_{i \geq 0} f_{i0}^{-1}(U_j)$
is an open subscheme of the limit we win.
\end{proof}


\begin{lemma}
\label{lemma-topology-limit}
In Situation \ref{situation-descent}.
\begin{enumerate}
\item We have $S_{set} = \lim_i S_{i, set}$ where $S_{set}$
indicates the underlying set of the scheme $S$.
\item We have $S_{top} = \lim_i S_{i, top}$ where $S_{top}$
indicates the underlying topological space of the scheme $S$.
\item If $s, s' \in S$ and $s'$ is not a specialization of $s$
then for some $i \in I$ the image $s'_i \in S_i$ of $s'$ is not
a specialization of the image $s_i \in S_i$ of $s$.
\item Add more easy facts on topology of $S$ here.
(Requirement: whatever is added should be easy in the affine case.)
\end{enumerate}
\end{lemma}

\begin{proof}
Part (1) is a special case of Lemma \href{https://stacks.math.columbia.edu/tag/0CUE}{lemma inverse limit sets (Tag 0CUE)}.

\medskip\noindent
Part (2) is a special case of Lemma \href{https://stacks.math.columbia.edu/tag/0CUF}{lemma inverse limit top (Tag 0CUF)}.

\medskip\noindent
Part (3) is a special case of Lemma \href{https://stacks.math.columbia.edu/tag/0CUG}{lemma inverse limit irreducibles (Tag 0CUG)}.
\end{proof}


\begin{lemma}
\label{lemma-scheme-over-limit}
Let $I$ be a directed set.
Let $(S_i, f_{ii'})$ be an inverse system of schemes over $I$.
Assume all the morphisms $f_{ii'} : S_i \to S_{i'}$ are affine,
Let $S = \lim_i S_i$. Let $0 \in I$.
Suppose that $T$ is a scheme over $S_0$.
Then
$$
T \times_{S_0} S = \lim_{i \geq 0} T \times_{S_0} S_i
$$
\end{lemma}

\begin{proof}
The right hand side is a scheme by
Lemma \ref{lemma-directed-inverse-system-has-limit}.
The equality is formal, see
Categories, Lemma \href{https://stacks.math.columbia.edu/tag/002M}{lemma colimits commute (Tag 002M)}.
\end{proof}


\begin{lemma}
\label{algebra-lemma-colimit-category-fp-algebras}
Suppose that $R = \colim_{\lambda \in \Lambda} R_\lambda$ is a directed colimit
of rings. Then the category of finitely presented $R$-algebras is
the colimit of the categories of finitely presented $R_\lambda$-algebras.
More precisely
\begin{enumerate}
\item Given a finitely presented $R$-algebra $A$ there exists a
$\lambda \in \Lambda$ and a finitely presented $R_\lambda$-algebra
$A_\lambda$ such that $A \cong A_\lambda \otimes_{R_\lambda} R$.
\item Given a $\lambda \in \Lambda$, finitely presented
$R_\lambda$-algebras $A_\lambda, B_\lambda$, and an $R$-algebra map
$\varphi : A_\lambda \otimes_{R_\lambda} R \to B_\lambda \otimes_{R_\lambda} R$,
then there exists a $\mu \geq \lambda$ and an $R_\mu$-algebra map
$\varphi_\mu : A_\lambda \otimes_{R_\lambda} R_\mu \to
B_\lambda \otimes_{R_\lambda} R_\mu$
such that $\varphi = \varphi_\mu \otimes 1_R$.
\item Given a $\lambda \in \Lambda$, finitely presented
$R_\lambda$-algebras $A_\lambda, B_\lambda$, and $R_\lambda$-algebra maps
$\varphi_\lambda, \psi_\lambda : A_\lambda \to B_\lambda$
such that $\varphi \otimes 1_R = \psi \otimes 1_R$, then
$\varphi \otimes 1_{R_\mu} = \psi \otimes 1_{R_\mu}$ for some
$\mu \geq \lambda$.
\end{enumerate}
\end{lemma}

\begin{proof}
To prove (1) choose a presentation
$A = R[x_1, \ldots, x_n]/(f_1, \ldots, f_m)$.
We can choose a $\lambda \in \Lambda$ and elements
$f_{\lambda, j} \in R_\lambda[x_1, \ldots, x_n]$ mapping to
$f_j \in R[x_1, \ldots, x_n]$.
Then we simply let
$A_\lambda =
R_\lambda[x_1, \ldots, x_n]/(f_{\lambda, 1}, \ldots, f_{\lambda, m})$.

\medskip\noindent
Parts (2) and (3) follow from
Lemma \ref{algebra-lemma-algebra-map-property-in-colimit}.
\end{proof}


\begin{lemma}
\label{algebra-lemma-algebra-map-property-in-colimit}
Let $A$ be a ring and let $B, C$ be $A$-algebras.
Suppose that $R = \colim_{i \in I} R_i$ is a directed colimit
of $A$-algebras.
\begin{enumerate}
\item If $B$ is a finite type $A$-algebra, and $u, u' : B \to C$ are
$A$-algebra maps such that
$u \otimes 1 = u' \otimes 1 : B \otimes_A R \to C \otimes_A R$
then for some $i$ we have
$u \otimes 1 = u' \otimes 1 : B \otimes_A R_i \to C \otimes_A R_i$.
\item If $C$ is a finite type $A$-algebra and $u : B \to C$ is an
$A$-algebra map such that
$u \otimes 1 : B \otimes_A R \to C \otimes_A R$ is surjective, then
for some $i$ the map $u \otimes 1 : B \otimes_A R_i \to C \otimes_A R_i$
is surjective.
\item If $C$ is of finite presentation over $A$ and
$v : C \otimes_A R \to B \otimes_A R$ is an $R$-algebra map, then there
exists an $i$ and an $R_i$-algebra map
$v_i : C \otimes_A R_i \to B \otimes_A R_i$ such that
$v = v_i \otimes 1$.
\item If $B$ is a finite type $A$-algebra, $C$ is a finitely presented
$A$-algebra, and
$u \otimes 1 : B \otimes_A R \to C \otimes_A R$ is an isomorphism, then
for some $i$ the map $u \otimes 1 : B \otimes_A R_i \to C \otimes_A R_i$
is an isomorphism.
\end{enumerate}
\end{lemma}

\begin{proof}
To prove (1) assume $u$ is as in (1) and
let $x_1, \ldots, x_m \in B$ be generators. Since
$B \otimes_A R = \colim_i B \otimes_A R_i$
we may pick an $i \in I$ such that $u(x_j) \otimes 1 = u'(x_j) \otimes 1$
in $B \otimes_A R_i$, $j = 1, \ldots, m$.
For such an $i$ we have
$u \otimes 1 = u' \otimes 1 : B \otimes_A R_i \to C \otimes_A R_i$.

\medskip\noindent
To prove (2) assume $u \otimes 1$ surjective and
let $y_1, \ldots, y_m \in C$ be generators. Since
$B \otimes_A R = \colim_i B \otimes_A R_i$
we may pick an $i \in I$ and $z_j \in B \otimes_A R_i$, $j = 1, \ldots, m$
whose images in $C \otimes_A R$ equal $y_j \otimes 1$.
For such an $i$ the map $u \otimes 1 : B \otimes_A R_i \to C \otimes_A R_i$
is surjective.

\medskip\noindent
To prove (3) let $c_1, \ldots, c_m \in C$ be generators. Let
$K = \Ker(A[x_1, \ldots, x_m] \to N)$ where the map is given by
the rule $x_j \mapsto \sum c_j$. Let $f_1, \ldots, f_t$
be generators for $K$ as an ideal in $A[x_1, \ldots, x_m]$.
We think of $f_j = f_j(x_1, \ldots, x_m)$ as a polynomial.
Since $B \otimes_A R = \colim_i B \otimes_A R_i$
we may pick an $i \in I$ and $z_j \in B \otimes_A R_i$, $j = 1, \ldots, m$
whose images in $B \otimes_A R$ equal $v(c_j \otimes 1)$.
We want to use the $z_j$ to define a map
$v_i : C \otimes_A R_i \to B \otimes_A R_i$.
Since $K \otimes_A R_i \to R_i[x_1, \ldots, x_m] \to C \otimes_A R_i \to 0$
is a presentation, it suffices to check that
$\xi_s = f_j(z_1, \ldots, z_m)$ is
zero in $B \otimes_A R_i$ for each $s = 1, \ldots, t$. This may not
be the case, but since the image of $\xi_s$ in $B \otimes_A R$ is zero
we see that it will be the case after increasing $i$ a bit.

\medskip\noindent
To prove (4) assume $u \otimes 1$ is an isomorphism, that
$B$ is a finite type $A$-algebra, and that $C$ is a finitely presented
$A$-algebra. Let $v : B \otimes_A R \to C \otimes_A R$ be an inverse to
$u \otimes 1$. Let $v_i : C \otimes_A R_i \to B \otimes_A R_i$ be as
in part (3). Apply part (1) to see that, after increasing $i$ we have
$v_i \circ (u \otimes 1) = \text{id}_{B \otimes_R R_i}$ and
$(u \otimes 1) \circ v_i = \text{id}_{C \otimes_R R_i}$.
\end{proof}






\begin{lemma}
\label{lemma-Noetherian-dvr-valuative-proper}
Let $S$ be a locally Noetherian scheme.
Let $f : X \to S$ be a morphism of finite type.
The following are equivalent:
\begin{enumerate}
\item The morphism $f$ is proper.
\item For any diagram (\ref{equation-valuative}) there exists exactly
one dotted arrow.
\item For all diagrams (\ref{equation-valuative}) with $A$ a discrete
valuation ring there exists exactly one dotted arrow.
\item For any irreducible component $X_0$ of $X$ with
generic point $\eta \in X_0$, for any discrete valuation ring
$A \subset K = \kappa(\eta)$ with fraction field $K$ and any
diagram (\ref{equation-valuative}) such that
the morphism $\Spec(K) \to X$ is the canonical one
(see
Schemes, Section \href{https://stacks.math.columbia.edu/tag/01J5}{section points (Tag 01J5)})
there exists exactly one dotted arrow.
\end{enumerate}
\end{lemma}

\begin{proof}
(1) implies (2) implies (3) implies (4).  We will now show (4) implies
(1).  As in the proof of Lemma \ref{lemma-Noetherian-dvr-valuative-separation},
we can reduce to the
case $S$ is affine, since properness is local on the base, and if $X
\to S$ satisfies (4), then $X_\alpha \to S_\alpha$ does as well for
open $S_\alpha \subset S$ and $X_\alpha = f^{-1}(S_\alpha)$.

\medskip\noindent
Now $S$ is a Noetherian scheme, and so $X$ is as well, since $X \to
S$ is of finite type.  Now we may use Chow's lemma
(Cohomology of Schemes, Lemma \ref{coherent-lemma-chow-Noetherian})
to get a surjective, proper, birational
$X' \to X$ and an immersion $X' \to \mathbf{P}^n_S$.  We wish to
show $X \to S$ is universally closed.  As in the proof of Lemma
\ref{lemma-limited-base-change}, it is enough to check that
$X' \to \mathbf{P}^n_S$ is a closed immersion.
For the sake of contradiction, assume that $X' \to
\mathbf{P}^n_S$ is not a closed immersion.  Then there is some $y
\in \mathbf{P}^n_S$ that is in the closure of the image of $X'$, but
is not in the image.  So $y$ is in the closure of the image of an
irreducible component $X_0'$ of $X'$, but not in the image.
Let $\bar X_0' \subset \mathbf{P}^n_S$ be the closure of
the image of $X_0'$. As $X' \to \mathbf{P}^n_S$ is an immersion
of Noetherian schemes, the morphism $X'_0 \to \bar X_0'$ is
open and dense. By
Algebra, Lemma \ref{algebra-lemma-exists-dvr}
or
Properties, Lemma \href{https://stacks.math.columbia.edu/tag/054F}{lemma locally Noetherian specialization dvr (Tag 054F)}
we can find a discrete valuation ring $A$ dominating
$\mathcal{O}_{\bar X_0', y}$ and with identical field
of fractions $K$. It is clear that
$K$ is the residue field at the generic point of $X_0'$.
Thus the solid commutative diagram
\begin{equation}
\label{equation-solid}
\xymatrix{
\Spec K \ar[r] \ar[d] & X' \ar [r] \ar[d] &
\mathbf{P}^n_S \ar[d] \\
\Spec A \ar@{-->}[r] \ar@{-->}[ru] \ar[urr] & X \ar[r] & S\\
}
\end{equation}
Note that the closed point of $A$ maps to $y \in \mathbf{P}^n_S$.  By
construction, there does not exist a set theoretic lift to $X'$.
As $X' \to X$ is birational, the image of $X'_0$ in $X$ is an
irreducible component $X_0$ of $X$ and $K$ is also identified with
the function field of $X_0$. Hence, as $X \to S$ is assumed to satisfy (4),
the dotted arrow $\Spec(A) \to X$ exists.
Since $X' \to X$ is proper, the dotted
arrow lifts to the dotted arrow $\Spec(A) \to X'$ (use Schemes,
Proposition \href{https://stacks.math.columbia.edu/tag/01KF}{proposition characterize universally closed (Tag 01KF)}).
We can compose this with the immersion $X' \to \mathbf{P}^n_S$ to obtain
another morphism (not depicted in the diagram) from
$\Spec(A) \to \mathbf{P}^n_S$.  Since $\mathbf{P}^n_S$
is proper over $S$, it satisfies (2), and so these two morphisms
agree.  This is a contradiction, for we have constructed the
forbidden lift of our original map $\Spec(A) \to \mathbf{P}^n_S$
to $X'$.
\end{proof}
\end{document}
